% ============================================================
% HSBI --- Quantitative Research Methods
% Mock Examination 1 (Lectures 1--4, ISLP Chapters 1--3)
% Prof. Dr. Christoph Weisser --- Summer Semester 2026
% ------------------------------------------------------------
% SINGLE SOURCE with solutions toggle:
%   Exam only:
%     pdflatex -jobname=Mock_Exam_1 mock_exam_1.tex
%   With solutions:
%     pdflatex -jobname=Mock_Exam_1_Solutions "\def\withsolutions{1}% ============================================================
% HSBI --- Quantitative Research Methods
% Mock Examination 1 (Lectures 1--4, ISLP Chapters 1--3)
% Prof. Dr. Christoph Weisser --- Summer Semester 2026
% ------------------------------------------------------------
% SINGLE SOURCE with solutions toggle:
%   Exam only:
%     pdflatex -jobname=Mock_Exam_1 mock_exam_1.tex
%   With solutions:
%     pdflatex -jobname=Mock_Exam_1_Solutions "\def\withsolutions{1}% ============================================================
% HSBI --- Quantitative Research Methods
% Mock Examination 1 (Lectures 1--4, ISLP Chapters 1--3)
% Prof. Dr. Christoph Weisser --- Summer Semester 2026
% ------------------------------------------------------------
% SINGLE SOURCE with solutions toggle:
%   Exam only:
%     pdflatex -jobname=Mock_Exam_1 mock_exam_1.tex
%   With solutions:
%     pdflatex -jobname=Mock_Exam_1_Solutions "\def\withsolutions{1}% ============================================================
% HSBI --- Quantitative Research Methods
% Mock Examination 1 (Lectures 1--4, ISLP Chapters 1--3)
% Prof. Dr. Christoph Weisser --- Summer Semester 2026
% ------------------------------------------------------------
% SINGLE SOURCE with solutions toggle:
%   Exam only:
%     pdflatex -jobname=Mock_Exam_1 mock_exam_1.tex
%   With solutions:
%     pdflatex -jobname=Mock_Exam_1_Solutions "\def\withsolutions{1}\input{mock_exam_1.tex}"
% ============================================================
\documentclass[11pt,a4paper]{article}
\usepackage[utf8]{inputenc}
\usepackage[T1]{fontenc}
\usepackage[english]{babel}
\usepackage[margin=2.2cm]{geometry}
\usepackage{amsmath,amssymb}
\usepackage{booktabs,array}
\usepackage{enumitem}
\usepackage{xcolor}
\usepackage{tcolorbox}
\tcbuselibrary{skins,breakable}

\setlength{\parindent}{0pt}

% Teal solution box (only shown when \withsolutions is defined)
\newtcolorbox{solutionbox}{enhanced,breakable,colback=teal!5,colframe=teal!55!black,
  boxrule=0pt,leftrule=2.5pt,arc=2pt,fonttitle=\bfseries,title=Solution,
  left=2mm,right=2mm,top=1mm,bottom=1.5mm}

\newcommand{\problem}[2]{\par\vspace{7mm}\noindent{\large\textbf{Problem #1}}%
  \hfill{\large\textbf{(#2 points)}}\par\nopagebreak\vspace{0.6mm}%
  \noindent\rule{\textwidth}{0.6pt}\par\nopagebreak\vspace{2.5mm}}
\newcommand{\pp}[1]{\textbf{[#1~P]}\enspace}

\begin{document}

% ------------------------------------------------------------
% Header block
% ------------------------------------------------------------
\begin{center}
  {\Large\textbf{HSBI --- Bielefeld University of Applied Sciences and Arts}}\\[1.2mm]
  {\large Quantitative Research Methods}\\[1mm]
  {\normalsize Summer Semester 2026 \quad$\cdot$\quad Examiner: Prof.\ Dr.\ Christoph Weisser}\\[4.5mm]
  {\LARGE\textbf{Mock Examination 1 (Lectures 1--4)}}\\[1.2mm]
  {\normalsize ISLP Chapters 1--3: Statistical Learning $\cdot$ $K$-Nearest Neighbours $\cdot$ Linear Regression}%
  \ifdefined\withsolutions\\[2.2mm]{\large\textcolor{teal!55!black}{\textbf{--- Version with detailed solutions ---}}}\fi
\end{center}

\vspace{2.5mm}
\begin{center}
\begin{tabular}{@{}ll@{}}
  \textbf{Duration:} & \textbf{90 minutes}\\
  \textbf{Total credit:} & \textbf{90 points}\\
  \textbf{Allowed aids:} & non-programmable calculator; one handwritten A4 sheet of notes (both sides)\\
\end{tabular}
\end{center}

\vspace{3mm}
Name: \rule[-1mm]{0.38\textwidth}{0.4pt}\hfill Matriculation no.: \rule[-1mm]{0.24\textwidth}{0.4pt}

\vspace{4.5mm}
\begin{center}
\renewcommand{\arraystretch}{1.35}
\begin{tabular}{|l|c|c|c|c|c|c|}
\hline
\textbf{Problem} & \textbf{P1} & \textbf{P2} & \textbf{P3} & \textbf{P4} & \textbf{P5} & \textbf{Total}\\\hline
Maximum points & 15 & 15 & 12 & 24 & 24 & \textbf{90}\\\hline
Achieved points & \hspace{9mm} & \hspace{9mm} & \hspace{9mm} & \hspace{9mm} & \hspace{9mm} & \hspace{11mm}\\\hline
\end{tabular}
\end{center}

\vspace{2.5mm}
\textbf{Instructions}
\begin{itemize}[nosep,leftmargin=5.5mm]
  \item Justify all answers. Numerical results without a derivation or a brief reasoning receive no credit.
  \item Round final numerical results to 3 decimal places unless stated otherwise. Small deviations caused by carrying rounded intermediate results are accepted.
  \item The number of points per problem roughly equals the recommended working time in minutes.
\end{itemize}
\vspace{1mm}\noindent\rule{\textwidth}{0.6pt}

% ============================================================
\problem{1: Fundamentals of statistical learning}{15}

\begin{enumerate}[label=(\alph*),leftmargin=7mm,itemsep=2.5mm]
  \item \pp{4} Explain the difference between using a statistical learning method for
    \emph{prediction} and for \emph{inference}. Give one concrete example of a question
    for each of the two goals.
  \item \pp{4} Contrast \emph{parametric} and \emph{non-parametric} methods for estimating
    the unknown function $f$ in $Y=f(X)+\varepsilon$: describe the two approaches briefly
    and state one advantage of each. Which general trade-off between \emph{flexibility}
    and \emph{interpretability} lies behind this choice?
  \item \pp{3} Classify each of the following scenarios as \emph{regression},
    \emph{classification} or \emph{unsupervised learning}, with a one-sentence justification:
    \begin{enumerate}[label=(\roman*),nosep,topsep=1mm]
      \item A telecom provider has monthly usage, contract length and age for 800 customers
        and wants to predict whether each customer will cancel the contract next month.
      \item An economist records profit, number of employees and industry for 90 firms in
        order to understand which of these factors are associated with CEO salary, and how strongly.
      \item An online retailer has the shopping baskets of 5{,}000 customers and wants to
        discover groups of customers with similar purchasing behaviour.
    \end{enumerate}
  \item \pp{4} Explain why the training MSE of a method is in general \emph{smaller} than its
    test MSE. What is \emph{overfitting}, and how does it show up when training and test MSE
    are plotted against model flexibility?
\end{enumerate}

\ifdefined\withsolutions
\begin{solutionbox}
\textbf{(a) Prediction vs.\ inference.}
\emph{Prediction:} we want an accurate output $\hat{Y}=\hat{f}(X)$ for new inputs $X$; the exact
form of $\hat{f}$ is secondary, it may be treated as a black box (accuracy is what counts).
Example: predict whether an incoming e-mail is spam, or predict next month's demand for a product.
\emph{Inference:} we want to understand \emph{how} $Y$ changes with $X_1,\dots,X_p$: which
predictors matter, in which direction and how strongly. Here $\hat{f}$ must be interpretable.
Example: which advertising channels increase sales, and by how much per euro spent?

\medskip
\textbf{(b) Parametric vs.\ non-parametric.}
\emph{Parametric:} a two-step approach: (1) assume a functional form for $f$, e.g.\ the linear
model $f(X)=\beta_0+\beta_1X_1+\dots+\beta_pX_p$; (2) estimate the finite set of parameters
(e.g.\ by least squares). The problem of estimating $f$ is reduced to estimating a few parameters.
Advantage: needs relatively few observations, is simple to estimate and easy to interpret.
\emph{Non-parametric:} no explicit functional form is assumed; $\hat{f}$ is constructed to follow
the data directly (e.g.\ $K$-nearest neighbours, splines). Advantage: can capture a much wider
range of possible shapes of $f$, since no potentially wrong form is imposed.
The underlying trade-off: more flexible methods can reduce bias but need many observations, risk
overfitting and are harder to interpret; if inference is the goal, more restrictive
(interpretable) models are usually preferred --- this is the \emph{flexibility--interpretability}
trade-off.

\medskip
\textbf{(c) Scenarios.}
(i) \emph{Classification} (supervised): the response ``cancels: yes/no'' is qualitative;
the goal is prediction.
(ii) \emph{Regression} (supervised): the response CEO salary is quantitative;
the goal is inference (which factors are associated with salary).
(iii) \emph{Unsupervised learning}: there is no response variable; the goal is to find
clusters of similar customers.

\medskip
\textbf{(d) Training vs.\ test MSE and overfitting.}
The method is fitted by making the \emph{training} error small, so $\hat{f}$ adapts to the
training observations --- including their random noise. The training MSE therefore
systematically underestimates the error on fresh data, and for a given data set
training MSE $\le$ test MSE in general. As flexibility increases, the training MSE decreases
monotonically, while the test MSE first decreases and then increases (U-shape).
\emph{Overfitting} means the method follows patterns in the training noise that do not exist
in the population; a less flexible method would achieve a \emph{smaller test} MSE. It shows up
as a widening gap: very small training MSE combined with a rising test MSE.
\end{solutionbox}
\fi

% ============================================================
\problem{2: Bias--variance decomposition}{15}

For the prediction of a response at a fixed test point $x_0$, three methods of increasing
flexibility were analysed. A simulation study yielded the following values
(all in squared units of $Y$); the noise variance is $\operatorname{Var}(\varepsilon)=\sigma^2=1.00$:

\begin{center}
\begin{tabular}{lccc}
\toprule
 & \textbf{Method A} & \textbf{Method B} & \textbf{Method C}\\
 & (linear) & (moderately flexible) & (very flexible)\\
\midrule
$[\operatorname{Bias}(\hat{f}(x_0))]^2$ & 2.25 & 0.16 & 0.01\\
$\operatorname{Var}(\hat{f}(x_0))$      & 0.25 & 1.44 & 3.60\\
\bottomrule
\end{tabular}
\end{center}

\begin{enumerate}[label=(\alph*),leftmargin=7mm,itemsep=2.5mm]
  \item \pp{4} State the decomposition of the expected test MSE
    $E\!\left[(y_0-\hat{f}(x_0))^2\right]$ into its three components and explain in one
    sentence each what the three components measure.
  \item \pp{5} Compute the expected test MSE at $x_0$ for each of the three methods.
    Which method should be chosen, and why?
  \item \pp{3} What is the smallest expected test MSE at $x_0$ that \emph{any} method
    (however good) could achieve here? Justify your answer.
  \item \pp{3} Describe how squared bias, variance, training MSE and expected test MSE
    typically behave as the flexibility of a method increases, and explain how the numbers
    in the table illustrate the characteristic U-shape of the test MSE.
\end{enumerate}

\ifdefined\withsolutions
\begin{solutionbox}
\textbf{(a) Decomposition.}
\[
E\big[(y_0-\hat{f}(x_0))^2\big]
=\operatorname{Var}(\hat{f}(x_0))
+\big[\operatorname{Bias}(\hat{f}(x_0))\big]^2
+\operatorname{Var}(\varepsilon).
\]
\emph{Variance:} how much $\hat{f}(x_0)$ would change if the method were fitted on a different
training set. \emph{Squared bias:} the systematic error made by approximating the true $f$
by the (possibly too simple) model. \emph{Irreducible error} $\sigma^2$: the noise variance
of $\varepsilon$, which no method can remove.

\medskip
\textbf{(b) Expected test MSE.}
\begin{align*}
\text{A: } & 2.25+0.25+1.00 = 3.50\\
\text{B: } & 0.16+1.44+1.00 = \mathbf{2.60}\\
\text{C: } & 0.01+3.60+1.00 = 4.61
\end{align*}
Choose \textbf{Method B}: it has the smallest expected test MSE. Its moderate flexibility buys a
large bias reduction (2.25 $\to$ 0.16) at an acceptable variance cost, whereas the very flexible
Method C pays more in variance ($+2.16$) than it gains in bias ($-0.15$).

\medskip
\textbf{(c) Lower bound.}
$\sigma^2=1.00$. Squared bias and variance are non-negative, so
$E[(y_0-\hat{f}(x_0))^2]\ge \operatorname{Var}(\varepsilon)=1.00$ for every method.
Even the true $f$ (zero bias, zero variance) cannot predict the noise term $\varepsilon_0$:
$1.00$ is the \emph{irreducible error}.

\medskip
\textbf{(d) Behaviour with increasing flexibility.}
Squared bias decreases (typically monotonically), variance increases, and the training MSE
decreases monotonically (a more flexible fit follows the training data ever more closely).
The expected test MSE is the sum bias$^2$ + variance + $\sigma^2$: it first falls, while the
bias reduction dominates, and then rises, when the variance increase dominates --- the U-shape.
In the table: $3.50 \to 2.60 \to 4.61$ for A $\to$ B $\to$ C; the minimum is attained at the
intermediate flexibility (B).
\end{solutionbox}
\fi

% ============================================================
\problem{3: $K$-nearest neighbours classification}{12}

An online shop wants to predict whether a visitor makes a purchase. For seven past visitors,
two standardised activity scores $X_1$ (pages viewed) and $X_2$ (minutes on site) and the
outcome $Y$ (purchase: Yes/No) were recorded:

\begin{center}
\begin{tabular}{cccc}
\toprule
Obs. & $X_1$ & $X_2$ & $Y$\\
\midrule
1 & 4 & 3 & Yes\\
2 & 2 & 2 & No\\
3 & 6 & 2 & No\\
4 & 4 & 5 & No\\
5 & 1 & 1 & Yes\\
6 & 7 & 4 & Yes\\
7 & 5 & 6 & No\\
\bottomrule
\end{tabular}
\end{center}

A new visitor has scores $x_0=(X_1,X_2)=(4,2)$.

\begin{enumerate}[label=(\alph*),leftmargin=7mm,itemsep=2.5mm]
  \item \pp{4} Compute the Euclidean distance between $x_0$ and each of the seven
    observations (3 decimals; you may rank by squared distances). How is the new visitor
    classified with $K=1$?
  \item \pp{4} How is the new visitor classified with $K=3$? State also the estimated
    conditional probability $\widehat{\Pr}(Y=\text{Yes}\mid X=x_0)$ that the $K=3$
    classifier uses.
  \item \pp{2} How does the choice of $K$ affect the bias and the variance of the KNN
    classifier (small $K$ vs.\ large $K$)? What is the \emph{training} error rate of the
    1-nearest-neighbour classifier?
  \item \pp{2} What is the \emph{Bayes classifier}, and why can it not be used directly in
    practice? In what sense is KNN related to it?
\end{enumerate}

\ifdefined\withsolutions
\begin{solutionbox}
\textbf{(a) Distances and $K=1$.} With $d(x_0,x_i)=\sqrt{(X_{1i}-4)^2+(X_{2i}-2)^2}$:
\begin{center}
\begin{tabular}{lccccccc}
\toprule
Obs. & 1 & 2 & 3 & 4 & 5 & 6 & 7\\
\midrule
$d^2$ & 1 & 4 & 4 & 9 & 10 & 13 & 17\\
$d$ & 1.000 & 2.000 & 2.000 & 3.000 & 3.162 & 3.606 & 4.123\\
$Y$ & Yes & No & No & No & Yes & Yes & No\\
\bottomrule
\end{tabular}
\end{center}
$K=1$: the single nearest neighbour is Obs.~1 ($d=1.000$) with label Yes
$\Rightarrow$ classify as \textbf{Yes}.

\medskip
\textbf{(b) $K=3$.} The three nearest neighbours are Obs.~1 ($d=1.000$, Yes),
Obs.~2 ($d=2.000$, No) and Obs.~3 ($d=2.000$, No). Estimated conditional probabilities:
$\widehat{\Pr}(\text{Yes}\mid x_0)=1/3\approx 0.333$ and
$\widehat{\Pr}(\text{No}\mid x_0)=2/3\approx 0.667$.
Majority vote $\Rightarrow$ classify as \textbf{No}. (Note that $K=1$ and $K=3$ disagree ---
the decision depends on the flexibility of the classifier.)

\medskip
\textbf{(c) Effect of $K$.} Small $K$ (e.g.\ $K=1$): very flexible, jagged decision boundary
$\Rightarrow$ \emph{low bias, high variance}. Large $K$: averaging over many neighbours gives a
smooth, almost linear boundary $\Rightarrow$ \emph{high bias, low variance}. The training error
rate of 1-NN is $0$ (ignoring ties): each training point is its own nearest neighbour --- which
says nothing about its test error.

\medskip
\textbf{(d) Bayes classifier.} It assigns each $x$ to the class $j$ with the largest \emph{true}
conditional probability $\Pr(Y=j\mid X=x)$; it minimises the expected test error rate, and its
error (the \emph{Bayes error rate}) is the lower bound for all classifiers. It is unattainable
in practice because the true conditional distribution of $Y$ given $X$ is unknown. KNN mimics
it by \emph{estimating} these conditional probabilities locally from the $K$ nearest training
observations and then classifying to the largest estimated probability.
\end{solutionbox}
\fi

% ============================================================
\problem{4: Simple linear regression by hand}{24}

A caf\'e chain investigates the effect of local radio advertising on sales.
For $n=5$ weeks it records the number of radio spots per week $x$ and the weekly sales $y$
(in 1{,}000 EUR):

\begin{center}
\begin{tabular}{lccccc}
\toprule
Week $i$ & 1 & 2 & 3 & 4 & 5\\
\midrule
Spots $x_i$ & 1 & 2 & 3 & 4 & 5\\
Sales $y_i$ & 6 & 7 & 11 & 13 & 18\\
\bottomrule
\end{tabular}
\end{center}

You may use without proof: $\sum x_i = 15$, $\sum y_i = 55$, $\sum x_i^2 = 55$,
$\sum x_i y_i = 195$. Consider the model $Y=\beta_0+\beta_1 X+\varepsilon$.

\begin{enumerate}[label=(\alph*),leftmargin=7mm,itemsep=2.5mm]
  \item \pp{6} Compute the least-squares estimates $\hat{\beta}_1$ and $\hat{\beta}_0$, and
    write down the fitted regression line. Interpret both estimates in the context of the
    problem (with units).
  \item \pp{4} Compute the fitted values $\hat{y}_i$ and the residuals $e_i$ for all five
    weeks, and verify that $\sum_{i=1}^{5} e_i = 0$.
  \item \pp{3} Compute the residual sum of squares (RSS) and the residual standard error
    (RSE), and interpret the RSE.
  \item \pp{6} Compute $\operatorname{SE}(\hat{\beta}_1)=\operatorname{RSE}/\sqrt{\sum_i (x_i-\bar{x})^2}$
    and test $H_0\!:\beta_1=0$ against $H_1\!:\beta_1\neq 0$ at level $\alpha=0.05$ using the
    $t$-statistic. Also give the 95\,\% confidence interval for $\beta_1$.
    \emph{Use $t_{0.975,\,3}=3.182$ (97.5\,\% quantile of the $t$-distribution with 3 degrees
    of freedom).} State your conclusion in one sentence.
  \item \pp{3} Compute the total sum of squares (TSS) and the coefficient of determination
    $R^2$, and interpret $R^2$.
  \item \pp{2} Predict the weekly sales for a week with $x=6$ radio spots. Name one reason
    why this particular prediction should be treated with caution.
\end{enumerate}

\ifdefined\withsolutions
\begin{solutionbox}
\textbf{(a) Coefficient estimates.} $\bar{x}=15/5=3$, $\bar{y}=55/5=11$.
\[
S_{xy}=\sum x_iy_i-n\bar{x}\bar{y}=195-5\cdot 3\cdot 11=30,\qquad
S_{xx}=\sum x_i^2-n\bar{x}^2=55-5\cdot 9=10.
\]
\[
\hat{\beta}_1=\frac{S_{xy}}{S_{xx}}=\frac{30}{10}=3.000,\qquad
\hat{\beta}_0=\bar{y}-\hat{\beta}_1\bar{x}=11-3\cdot 3=2.000.
\]
Fitted line: $\hat{y}=2.000+3.000\,x$.
\emph{Interpretation:} each additional radio spot per week is associated with an expected
increase in weekly sales of $3.000$ thousand EUR ($=3{,}000$ EUR). The intercept $2.000$ is the
expected weekly sales ($2{,}000$ EUR) in a week without any spots --- a mild extrapolation,
since $x=0$ was not observed ($x$ ranges from 1 to 5).

\medskip
\textbf{(b) Fitted values and residuals.} $\hat{y}_i=2+3x_i$ and $e_i=y_i-\hat{y}_i$:
\begin{center}
\begin{tabular}{lccccc}
\toprule
$i$ & 1 & 2 & 3 & 4 & 5\\
\midrule
$\hat{y}_i$ & 5 & 8 & 11 & 14 & 17\\
$e_i$ & $1$ & $-1$ & $0$ & $-1$ & $1$\\
\bottomrule
\end{tabular}
\end{center}
$\sum_i e_i = 1-1+0-1+1 = 0$ \checkmark\ (this always holds for least squares with an intercept).

\medskip
\textbf{(c) RSS and RSE.}
$\operatorname{RSS}=\sum_i e_i^2 = 1+1+0+1+1 = 4.000$;
$\operatorname{RSE}=\sqrt{\operatorname{RSS}/(n-2)}=\sqrt{4/3}=1.155$ (thousand EUR).
\emph{Interpretation:} even if the true line were known, observed weekly sales would deviate
from it by about $1.155$ thousand EUR ($\approx 1{,}155$ EUR) on average --- an estimate of
$\sigma$, the standard deviation of $\varepsilon$.

\medskip
\textbf{(d) Standard error, $t$-test and confidence interval.}
\[
\operatorname{SE}(\hat{\beta}_1)=\frac{1.1547}{\sqrt{10}}=\frac{1.1547}{3.1623}=0.365,\qquad
t=\frac{\hat{\beta}_1-0}{\operatorname{SE}(\hat{\beta}_1)}=\frac{3.000}{0.3651}=8.216
\quad(\text{df}=n-2=3).
\]
Since $|t|=8.216 > 3.182=t_{0.975,3}$, reject $H_0$ at the 5\,\% level.
95\,\% confidence interval:
\[
3.000 \pm 3.182\cdot 0.3651 = 3.000 \pm 1.162 = [1.838,\; 4.162].
\]
Zero is not contained in the interval --- consistent with the test.
\emph{Conclusion:} there is a statistically significant (positive) linear association between
the number of radio spots and weekly sales.

\medskip
\textbf{(e) TSS and $R^2$.}
$\operatorname{TSS}=\sum_i (y_i-\bar{y})^2 = 25+16+0+4+49 = 94.000$;
\[
R^2=1-\frac{\operatorname{RSS}}{\operatorname{TSS}}=1-\frac{4}{94}=\frac{90}{94}=0.957.
\]
About $95.7\,\%$ of the variance of weekly sales is explained by the linear regression on the
number of radio spots. (In simple regression $R^2=r^2$, so $r=\sqrt{0.957}\approx 0.979$.)

\medskip
\textbf{(f) Prediction at $x=6$.}
$\hat{y}(6)=2+3\cdot 6=20.000$ thousand EUR ($=20{,}000$ EUR).
Caution: $x=6$ lies \emph{outside} the observed range $x\in[1,5]$ --- an extrapolation; the
linear relationship has not been observed there and may not hold. (Also, $n=5$ is very small.)
\end{solutionbox}
\fi

% ============================================================
\problem{5: Multiple linear regression --- interpreting Python output}{24}

A consultancy analyses the salaries of $n=60$ employees of a software company.
The variables are \texttt{salary} (annual salary in 1{,}000 EUR), \texttt{experience}
(professional experience in years) and the dummy variable \texttt{female}
($=1$ for female employees, $=0$ for male employees). The model
\[
\texttt{salary}=\beta_0+\beta_1\,\texttt{experience}+\beta_2\,\texttt{female}
+\beta_3\,(\texttt{experience}\times\texttt{female})+\varepsilon
\]
was fitted by least squares in Python (\texttt{statsmodels}), with the following output:

\begin{center}
\begin{tabular}{lrrrr}
\toprule
 & coef & std err & $t$ & $P>|t|$\\
\midrule
Intercept                    & 46.2000 & 1.100 & 42.000 & $<$0.0001\\
experience                   &  2.3500 & 0.250 &  9.400 & $<$0.0001\\
female                       & $-3.1000$ & 1.550 & $-2.000$ & 0.0504\\
experience:female            & $-0.8500$ & 0.350 & $-2.429$ & 0.0184\\
\bottomrule
\end{tabular}

\medskip
$n=60$; \quad $R^2 = 0.712$; \quad adjusted $R^2 = 0.697$;\\[0.5mm]
$F$-statistic $= 46.15$ on 3 and 56 degrees of freedom, \ Prob($F$-statistic) $< 0.0001$.
\end{center}

\begin{enumerate}[label=(\alph*),leftmargin=7mm,itemsep=2.5mm]
  \item \pp{6} Interpret all four estimated coefficients precisely in the context of the
    problem (with units), paying particular attention to the dummy variable and the
    interaction term.
  \item \pp{4} Which coefficients are significantly different from zero at level
    $\alpha=0.05$? Verify for the interaction term that the reported $t$-value is consistent
    with the reported coefficient and standard error. Should \texttt{female} be removed from
    the model because of its $p$-value? Discuss briefly.
  \item \pp{4} Write down the fitted regression equation separately for men and for women.
  \item \pp{4} Predict the annual salary of (i) a man with 10 years of experience and (ii) a
    woman with 10 years of experience. What is the estimated salary gap between women and men
    at 10 years of experience, and how can it be written in terms of $\hat{\beta}_2$ and
    $\hat{\beta}_3$?
  \item \pp{3} Interpret $R^2$. State the null hypothesis of the reported $F$-test and the
    conclusion that follows from the output.
  \item \pp{3} Model diagnostics:
    \begin{enumerate}[label=(\roman*),nosep,topsep=1mm]
      \item A plot of the residuals against the fitted values shows a clear U-shape
        (positive residuals for small and large fitted values, negative in between).
        What does this indicate, and name one possible remedy.
      \item In an extended model, the predictor \texttt{age} has a variance inflation factor
        of $\operatorname{VIF}=12$. What does this mean, and what is the practical consequence
        for the regression results?
    \end{enumerate}
\end{enumerate}

\ifdefined\withsolutions
\begin{solutionbox}
\textbf{(a) Interpretation of the coefficients} (salary in 1{,}000 EUR per year):
\begin{itemize}[nosep,leftmargin=5mm]
  \item \emph{Intercept} $46.20$: expected annual salary of a \emph{man}
    (\texttt{female}=0) with 0 years of experience: $46{,}200$ EUR.
  \item \emph{experience} $2.35$: for \emph{men}, each additional year of experience increases
    the expected salary by $2{,}350$ EUR. (Because of the interaction, this slope refers to the
    baseline group \texttt{female}=0 only.)
  \item \emph{female} $-3.10$: at 0 years of experience, women are expected to earn
    $3{,}100$ EUR \emph{less} than men --- the difference between the two intercepts.
  \item \emph{experience:female} $-0.85$: the experience slope for women is $0.85$ \emph{lower}
    than for men: women gain $2.35-0.85=1.50$ thousand EUR per additional year. Equivalently,
    the expected salary gap between men and women \emph{widens} by $850$ EUR with every year
    of experience.
\end{itemize}

\medskip
\textbf{(b) Significance at $\alpha=0.05$.}
Intercept and \texttt{experience}: $p<0.0001$ $\Rightarrow$ highly significant.
\texttt{experience:female}: $p=0.0184<0.05$ $\Rightarrow$ significant.
Consistency check: $t=\hat{\beta}_3/\operatorname{SE}(\hat{\beta}_3)=-0.85/0.35=-2.429$
\checkmark\ (matches the output; $p$-value from $t_{56}$).
\texttt{female}: $p=0.0504>0.05$ $\Rightarrow$ \emph{not} significant at exactly the 5\,\%
level, though only marginally. It should nevertheless \emph{not} be removed: by the
\emph{hierarchy principle}, if an interaction ($\texttt{experience}\times\texttt{female}$) is
kept in the model, the corresponding main effects should be kept as well --- otherwise the
interaction changes its meaning. Moreover, $p=0.0504$ is borderline evidence, not evidence of
no effect.

\medskip
\textbf{(c) Fitted equations per group.}
\begin{align*}
\text{Men } (\texttt{female}=0):\quad & \widehat{\texttt{salary}} = 46.20 + 2.35\,\texttt{experience}\\
\text{Women } (\texttt{female}=1):\quad & \widehat{\texttt{salary}} = (46.20-3.10) + (2.35-0.85)\,\texttt{experience}\\
 & \phantom{\widehat{\texttt{salary}}} = 43.10 + 1.50\,\texttt{experience}
\end{align*}
Two different intercepts \emph{and} two different slopes (that is exactly what the interaction allows).

\medskip
\textbf{(d) Predictions at 10 years of experience.}
\begin{align*}
\text{Man: } & 46.20 + 2.35\cdot 10 = 69.70 \;\Rightarrow\; 69{,}700 \text{ EUR}\\
\text{Woman: } & 43.10 + 1.50\cdot 10 = 58.10 \;\Rightarrow\; 58{,}100 \text{ EUR}
\end{align*}
Estimated gap: $58.10-69.70=-11.60$, i.e.\ women earn $11{,}600$ EUR less. In terms of the
coefficients: $\hat{\beta}_2+10\,\hat{\beta}_3=-3.10-8.50=-11.60$ thousand EUR.

\medskip
\textbf{(e) $R^2$ and $F$-test.}
$R^2=0.712$: the three model terms explain $71.2\,\%$ of the variance of salaries in the sample.
$F$-test: $H_0\!: \beta_1=\beta_2=\beta_3=0$ (none of the predictors is related to salary)
against the alternative that at least one $\beta_j\neq 0$. With $F=46.15$ (3 and 56 df) and
$p<0.0001$, $H_0$ is clearly rejected: the model is useful overall.

\medskip
\textbf{(f) Diagnostics.}
(i) A U-shaped residual-vs-fitted plot indicates \emph{non-linearity}: the linear model misses
systematic curvature in the relationship. Remedy: include a non-linear term, e.g.\ add
$\texttt{experience}^2$ (polynomial regression), or transform variables
(e.g.\ $\log(\texttt{salary})$).
(ii) $\operatorname{VIF}_j = 1/(1-R^2_{X_j\mid X_{-j}})$; $\operatorname{VIF}=12$ means that
regressing \texttt{age} on the other predictors gives $R^2\approx 0.917$: \texttt{age} is
almost a linear combination of the other predictors (strong \emph{multicollinearity}; common
rules of thumb flag VIF $>5$ or $>10$). Consequence: $\operatorname{SE}(\hat{\beta}_{\texttt{age}})$
is inflated by a factor $\sqrt{12}\approx 3.5$ relative to the uncorrelated case --- the
coefficient is estimated very imprecisely, its sign may flip between samples, and its $t$-test
loses power, even though the model's overall fit and predictions may be unaffected.
\end{solutionbox}
\fi

\vspace{6mm}
\begin{center}
\rule{0.35\textwidth}{0.4pt}\\[1mm]
\textit{End of the examination --- good luck!}
\end{center}

\end{document}
"
% ============================================================
\documentclass[11pt,a4paper]{article}
\usepackage[utf8]{inputenc}
\usepackage[T1]{fontenc}
\usepackage[english]{babel}
\usepackage[margin=2.2cm]{geometry}
\usepackage{amsmath,amssymb}
\usepackage{booktabs,array}
\usepackage{enumitem}
\usepackage{xcolor}
\usepackage{tcolorbox}
\tcbuselibrary{skins,breakable}

\setlength{\parindent}{0pt}

% Teal solution box (only shown when \withsolutions is defined)
\newtcolorbox{solutionbox}{enhanced,breakable,colback=teal!5,colframe=teal!55!black,
  boxrule=0pt,leftrule=2.5pt,arc=2pt,fonttitle=\bfseries,title=Solution,
  left=2mm,right=2mm,top=1mm,bottom=1.5mm}

\newcommand{\problem}[2]{\par\vspace{7mm}\noindent{\large\textbf{Problem #1}}%
  \hfill{\large\textbf{(#2 points)}}\par\nopagebreak\vspace{0.6mm}%
  \noindent\rule{\textwidth}{0.6pt}\par\nopagebreak\vspace{2.5mm}}
\newcommand{\pp}[1]{\textbf{[#1~P]}\enspace}

\begin{document}

% ------------------------------------------------------------
% Header block
% ------------------------------------------------------------
\begin{center}
  {\Large\textbf{HSBI --- Bielefeld University of Applied Sciences and Arts}}\\[1.2mm]
  {\large Quantitative Research Methods}\\[1mm]
  {\normalsize Summer Semester 2026 \quad$\cdot$\quad Examiner: Prof.\ Dr.\ Christoph Weisser}\\[4.5mm]
  {\LARGE\textbf{Mock Examination 1 (Lectures 1--4)}}\\[1.2mm]
  {\normalsize ISLP Chapters 1--3: Statistical Learning $\cdot$ $K$-Nearest Neighbours $\cdot$ Linear Regression}%
  \ifdefined\withsolutions\\[2.2mm]{\large\textcolor{teal!55!black}{\textbf{--- Version with detailed solutions ---}}}\fi
\end{center}

\vspace{2.5mm}
\begin{center}
\begin{tabular}{@{}ll@{}}
  \textbf{Duration:} & \textbf{90 minutes}\\
  \textbf{Total credit:} & \textbf{90 points}\\
  \textbf{Allowed aids:} & non-programmable calculator; one handwritten A4 sheet of notes (both sides)\\
\end{tabular}
\end{center}

\vspace{3mm}
Name: \rule[-1mm]{0.38\textwidth}{0.4pt}\hfill Matriculation no.: \rule[-1mm]{0.24\textwidth}{0.4pt}

\vspace{4.5mm}
\begin{center}
\renewcommand{\arraystretch}{1.35}
\begin{tabular}{|l|c|c|c|c|c|c|}
\hline
\textbf{Problem} & \textbf{P1} & \textbf{P2} & \textbf{P3} & \textbf{P4} & \textbf{P5} & \textbf{Total}\\\hline
Maximum points & 15 & 15 & 12 & 24 & 24 & \textbf{90}\\\hline
Achieved points & \hspace{9mm} & \hspace{9mm} & \hspace{9mm} & \hspace{9mm} & \hspace{9mm} & \hspace{11mm}\\\hline
\end{tabular}
\end{center}

\vspace{2.5mm}
\textbf{Instructions}
\begin{itemize}[nosep,leftmargin=5.5mm]
  \item Justify all answers. Numerical results without a derivation or a brief reasoning receive no credit.
  \item Round final numerical results to 3 decimal places unless stated otherwise. Small deviations caused by carrying rounded intermediate results are accepted.
  \item The number of points per problem roughly equals the recommended working time in minutes.
\end{itemize}
\vspace{1mm}\noindent\rule{\textwidth}{0.6pt}

% ============================================================
\problem{1: Fundamentals of statistical learning}{15}

\begin{enumerate}[label=(\alph*),leftmargin=7mm,itemsep=2.5mm]
  \item \pp{4} Explain the difference between using a statistical learning method for
    \emph{prediction} and for \emph{inference}. Give one concrete example of a question
    for each of the two goals.
  \item \pp{4} Contrast \emph{parametric} and \emph{non-parametric} methods for estimating
    the unknown function $f$ in $Y=f(X)+\varepsilon$: describe the two approaches briefly
    and state one advantage of each. Which general trade-off between \emph{flexibility}
    and \emph{interpretability} lies behind this choice?
  \item \pp{3} Classify each of the following scenarios as \emph{regression},
    \emph{classification} or \emph{unsupervised learning}, with a one-sentence justification:
    \begin{enumerate}[label=(\roman*),nosep,topsep=1mm]
      \item A telecom provider has monthly usage, contract length and age for 800 customers
        and wants to predict whether each customer will cancel the contract next month.
      \item An economist records profit, number of employees and industry for 90 firms in
        order to understand which of these factors are associated with CEO salary, and how strongly.
      \item An online retailer has the shopping baskets of 5{,}000 customers and wants to
        discover groups of customers with similar purchasing behaviour.
    \end{enumerate}
  \item \pp{4} Explain why the training MSE of a method is in general \emph{smaller} than its
    test MSE. What is \emph{overfitting}, and how does it show up when training and test MSE
    are plotted against model flexibility?
\end{enumerate}

\ifdefined\withsolutions
\begin{solutionbox}
\textbf{(a) Prediction vs.\ inference.}
\emph{Prediction:} we want an accurate output $\hat{Y}=\hat{f}(X)$ for new inputs $X$; the exact
form of $\hat{f}$ is secondary, it may be treated as a black box (accuracy is what counts).
Example: predict whether an incoming e-mail is spam, or predict next month's demand for a product.
\emph{Inference:} we want to understand \emph{how} $Y$ changes with $X_1,\dots,X_p$: which
predictors matter, in which direction and how strongly. Here $\hat{f}$ must be interpretable.
Example: which advertising channels increase sales, and by how much per euro spent?

\medskip
\textbf{(b) Parametric vs.\ non-parametric.}
\emph{Parametric:} a two-step approach: (1) assume a functional form for $f$, e.g.\ the linear
model $f(X)=\beta_0+\beta_1X_1+\dots+\beta_pX_p$; (2) estimate the finite set of parameters
(e.g.\ by least squares). The problem of estimating $f$ is reduced to estimating a few parameters.
Advantage: needs relatively few observations, is simple to estimate and easy to interpret.
\emph{Non-parametric:} no explicit functional form is assumed; $\hat{f}$ is constructed to follow
the data directly (e.g.\ $K$-nearest neighbours, splines). Advantage: can capture a much wider
range of possible shapes of $f$, since no potentially wrong form is imposed.
The underlying trade-off: more flexible methods can reduce bias but need many observations, risk
overfitting and are harder to interpret; if inference is the goal, more restrictive
(interpretable) models are usually preferred --- this is the \emph{flexibility--interpretability}
trade-off.

\medskip
\textbf{(c) Scenarios.}
(i) \emph{Classification} (supervised): the response ``cancels: yes/no'' is qualitative;
the goal is prediction.
(ii) \emph{Regression} (supervised): the response CEO salary is quantitative;
the goal is inference (which factors are associated with salary).
(iii) \emph{Unsupervised learning}: there is no response variable; the goal is to find
clusters of similar customers.

\medskip
\textbf{(d) Training vs.\ test MSE and overfitting.}
The method is fitted by making the \emph{training} error small, so $\hat{f}$ adapts to the
training observations --- including their random noise. The training MSE therefore
systematically underestimates the error on fresh data, and for a given data set
training MSE $\le$ test MSE in general. As flexibility increases, the training MSE decreases
monotonically, while the test MSE first decreases and then increases (U-shape).
\emph{Overfitting} means the method follows patterns in the training noise that do not exist
in the population; a less flexible method would achieve a \emph{smaller test} MSE. It shows up
as a widening gap: very small training MSE combined with a rising test MSE.
\end{solutionbox}
\fi

% ============================================================
\problem{2: Bias--variance decomposition}{15}

For the prediction of a response at a fixed test point $x_0$, three methods of increasing
flexibility were analysed. A simulation study yielded the following values
(all in squared units of $Y$); the noise variance is $\operatorname{Var}(\varepsilon)=\sigma^2=1.00$:

\begin{center}
\begin{tabular}{lccc}
\toprule
 & \textbf{Method A} & \textbf{Method B} & \textbf{Method C}\\
 & (linear) & (moderately flexible) & (very flexible)\\
\midrule
$[\operatorname{Bias}(\hat{f}(x_0))]^2$ & 2.25 & 0.16 & 0.01\\
$\operatorname{Var}(\hat{f}(x_0))$      & 0.25 & 1.44 & 3.60\\
\bottomrule
\end{tabular}
\end{center}

\begin{enumerate}[label=(\alph*),leftmargin=7mm,itemsep=2.5mm]
  \item \pp{4} State the decomposition of the expected test MSE
    $E\!\left[(y_0-\hat{f}(x_0))^2\right]$ into its three components and explain in one
    sentence each what the three components measure.
  \item \pp{5} Compute the expected test MSE at $x_0$ for each of the three methods.
    Which method should be chosen, and why?
  \item \pp{3} What is the smallest expected test MSE at $x_0$ that \emph{any} method
    (however good) could achieve here? Justify your answer.
  \item \pp{3} Describe how squared bias, variance, training MSE and expected test MSE
    typically behave as the flexibility of a method increases, and explain how the numbers
    in the table illustrate the characteristic U-shape of the test MSE.
\end{enumerate}

\ifdefined\withsolutions
\begin{solutionbox}
\textbf{(a) Decomposition.}
\[
E\big[(y_0-\hat{f}(x_0))^2\big]
=\operatorname{Var}(\hat{f}(x_0))
+\big[\operatorname{Bias}(\hat{f}(x_0))\big]^2
+\operatorname{Var}(\varepsilon).
\]
\emph{Variance:} how much $\hat{f}(x_0)$ would change if the method were fitted on a different
training set. \emph{Squared bias:} the systematic error made by approximating the true $f$
by the (possibly too simple) model. \emph{Irreducible error} $\sigma^2$: the noise variance
of $\varepsilon$, which no method can remove.

\medskip
\textbf{(b) Expected test MSE.}
\begin{align*}
\text{A: } & 2.25+0.25+1.00 = 3.50\\
\text{B: } & 0.16+1.44+1.00 = \mathbf{2.60}\\
\text{C: } & 0.01+3.60+1.00 = 4.61
\end{align*}
Choose \textbf{Method B}: it has the smallest expected test MSE. Its moderate flexibility buys a
large bias reduction (2.25 $\to$ 0.16) at an acceptable variance cost, whereas the very flexible
Method C pays more in variance ($+2.16$) than it gains in bias ($-0.15$).

\medskip
\textbf{(c) Lower bound.}
$\sigma^2=1.00$. Squared bias and variance are non-negative, so
$E[(y_0-\hat{f}(x_0))^2]\ge \operatorname{Var}(\varepsilon)=1.00$ for every method.
Even the true $f$ (zero bias, zero variance) cannot predict the noise term $\varepsilon_0$:
$1.00$ is the \emph{irreducible error}.

\medskip
\textbf{(d) Behaviour with increasing flexibility.}
Squared bias decreases (typically monotonically), variance increases, and the training MSE
decreases monotonically (a more flexible fit follows the training data ever more closely).
The expected test MSE is the sum bias$^2$ + variance + $\sigma^2$: it first falls, while the
bias reduction dominates, and then rises, when the variance increase dominates --- the U-shape.
In the table: $3.50 \to 2.60 \to 4.61$ for A $\to$ B $\to$ C; the minimum is attained at the
intermediate flexibility (B).
\end{solutionbox}
\fi

% ============================================================
\problem{3: $K$-nearest neighbours classification}{12}

An online shop wants to predict whether a visitor makes a purchase. For seven past visitors,
two standardised activity scores $X_1$ (pages viewed) and $X_2$ (minutes on site) and the
outcome $Y$ (purchase: Yes/No) were recorded:

\begin{center}
\begin{tabular}{cccc}
\toprule
Obs. & $X_1$ & $X_2$ & $Y$\\
\midrule
1 & 4 & 3 & Yes\\
2 & 2 & 2 & No\\
3 & 6 & 2 & No\\
4 & 4 & 5 & No\\
5 & 1 & 1 & Yes\\
6 & 7 & 4 & Yes\\
7 & 5 & 6 & No\\
\bottomrule
\end{tabular}
\end{center}

A new visitor has scores $x_0=(X_1,X_2)=(4,2)$.

\begin{enumerate}[label=(\alph*),leftmargin=7mm,itemsep=2.5mm]
  \item \pp{4} Compute the Euclidean distance between $x_0$ and each of the seven
    observations (3 decimals; you may rank by squared distances). How is the new visitor
    classified with $K=1$?
  \item \pp{4} How is the new visitor classified with $K=3$? State also the estimated
    conditional probability $\widehat{\Pr}(Y=\text{Yes}\mid X=x_0)$ that the $K=3$
    classifier uses.
  \item \pp{2} How does the choice of $K$ affect the bias and the variance of the KNN
    classifier (small $K$ vs.\ large $K$)? What is the \emph{training} error rate of the
    1-nearest-neighbour classifier?
  \item \pp{2} What is the \emph{Bayes classifier}, and why can it not be used directly in
    practice? In what sense is KNN related to it?
\end{enumerate}

\ifdefined\withsolutions
\begin{solutionbox}
\textbf{(a) Distances and $K=1$.} With $d(x_0,x_i)=\sqrt{(X_{1i}-4)^2+(X_{2i}-2)^2}$:
\begin{center}
\begin{tabular}{lccccccc}
\toprule
Obs. & 1 & 2 & 3 & 4 & 5 & 6 & 7\\
\midrule
$d^2$ & 1 & 4 & 4 & 9 & 10 & 13 & 17\\
$d$ & 1.000 & 2.000 & 2.000 & 3.000 & 3.162 & 3.606 & 4.123\\
$Y$ & Yes & No & No & No & Yes & Yes & No\\
\bottomrule
\end{tabular}
\end{center}
$K=1$: the single nearest neighbour is Obs.~1 ($d=1.000$) with label Yes
$\Rightarrow$ classify as \textbf{Yes}.

\medskip
\textbf{(b) $K=3$.} The three nearest neighbours are Obs.~1 ($d=1.000$, Yes),
Obs.~2 ($d=2.000$, No) and Obs.~3 ($d=2.000$, No). Estimated conditional probabilities:
$\widehat{\Pr}(\text{Yes}\mid x_0)=1/3\approx 0.333$ and
$\widehat{\Pr}(\text{No}\mid x_0)=2/3\approx 0.667$.
Majority vote $\Rightarrow$ classify as \textbf{No}. (Note that $K=1$ and $K=3$ disagree ---
the decision depends on the flexibility of the classifier.)

\medskip
\textbf{(c) Effect of $K$.} Small $K$ (e.g.\ $K=1$): very flexible, jagged decision boundary
$\Rightarrow$ \emph{low bias, high variance}. Large $K$: averaging over many neighbours gives a
smooth, almost linear boundary $\Rightarrow$ \emph{high bias, low variance}. The training error
rate of 1-NN is $0$ (ignoring ties): each training point is its own nearest neighbour --- which
says nothing about its test error.

\medskip
\textbf{(d) Bayes classifier.} It assigns each $x$ to the class $j$ with the largest \emph{true}
conditional probability $\Pr(Y=j\mid X=x)$; it minimises the expected test error rate, and its
error (the \emph{Bayes error rate}) is the lower bound for all classifiers. It is unattainable
in practice because the true conditional distribution of $Y$ given $X$ is unknown. KNN mimics
it by \emph{estimating} these conditional probabilities locally from the $K$ nearest training
observations and then classifying to the largest estimated probability.
\end{solutionbox}
\fi

% ============================================================
\problem{4: Simple linear regression by hand}{24}

A caf\'e chain investigates the effect of local radio advertising on sales.
For $n=5$ weeks it records the number of radio spots per week $x$ and the weekly sales $y$
(in 1{,}000 EUR):

\begin{center}
\begin{tabular}{lccccc}
\toprule
Week $i$ & 1 & 2 & 3 & 4 & 5\\
\midrule
Spots $x_i$ & 1 & 2 & 3 & 4 & 5\\
Sales $y_i$ & 6 & 7 & 11 & 13 & 18\\
\bottomrule
\end{tabular}
\end{center}

You may use without proof: $\sum x_i = 15$, $\sum y_i = 55$, $\sum x_i^2 = 55$,
$\sum x_i y_i = 195$. Consider the model $Y=\beta_0+\beta_1 X+\varepsilon$.

\begin{enumerate}[label=(\alph*),leftmargin=7mm,itemsep=2.5mm]
  \item \pp{6} Compute the least-squares estimates $\hat{\beta}_1$ and $\hat{\beta}_0$, and
    write down the fitted regression line. Interpret both estimates in the context of the
    problem (with units).
  \item \pp{4} Compute the fitted values $\hat{y}_i$ and the residuals $e_i$ for all five
    weeks, and verify that $\sum_{i=1}^{5} e_i = 0$.
  \item \pp{3} Compute the residual sum of squares (RSS) and the residual standard error
    (RSE), and interpret the RSE.
  \item \pp{6} Compute $\operatorname{SE}(\hat{\beta}_1)=\operatorname{RSE}/\sqrt{\sum_i (x_i-\bar{x})^2}$
    and test $H_0\!:\beta_1=0$ against $H_1\!:\beta_1\neq 0$ at level $\alpha=0.05$ using the
    $t$-statistic. Also give the 95\,\% confidence interval for $\beta_1$.
    \emph{Use $t_{0.975,\,3}=3.182$ (97.5\,\% quantile of the $t$-distribution with 3 degrees
    of freedom).} State your conclusion in one sentence.
  \item \pp{3} Compute the total sum of squares (TSS) and the coefficient of determination
    $R^2$, and interpret $R^2$.
  \item \pp{2} Predict the weekly sales for a week with $x=6$ radio spots. Name one reason
    why this particular prediction should be treated with caution.
\end{enumerate}

\ifdefined\withsolutions
\begin{solutionbox}
\textbf{(a) Coefficient estimates.} $\bar{x}=15/5=3$, $\bar{y}=55/5=11$.
\[
S_{xy}=\sum x_iy_i-n\bar{x}\bar{y}=195-5\cdot 3\cdot 11=30,\qquad
S_{xx}=\sum x_i^2-n\bar{x}^2=55-5\cdot 9=10.
\]
\[
\hat{\beta}_1=\frac{S_{xy}}{S_{xx}}=\frac{30}{10}=3.000,\qquad
\hat{\beta}_0=\bar{y}-\hat{\beta}_1\bar{x}=11-3\cdot 3=2.000.
\]
Fitted line: $\hat{y}=2.000+3.000\,x$.
\emph{Interpretation:} each additional radio spot per week is associated with an expected
increase in weekly sales of $3.000$ thousand EUR ($=3{,}000$ EUR). The intercept $2.000$ is the
expected weekly sales ($2{,}000$ EUR) in a week without any spots --- a mild extrapolation,
since $x=0$ was not observed ($x$ ranges from 1 to 5).

\medskip
\textbf{(b) Fitted values and residuals.} $\hat{y}_i=2+3x_i$ and $e_i=y_i-\hat{y}_i$:
\begin{center}
\begin{tabular}{lccccc}
\toprule
$i$ & 1 & 2 & 3 & 4 & 5\\
\midrule
$\hat{y}_i$ & 5 & 8 & 11 & 14 & 17\\
$e_i$ & $1$ & $-1$ & $0$ & $-1$ & $1$\\
\bottomrule
\end{tabular}
\end{center}
$\sum_i e_i = 1-1+0-1+1 = 0$ \checkmark\ (this always holds for least squares with an intercept).

\medskip
\textbf{(c) RSS and RSE.}
$\operatorname{RSS}=\sum_i e_i^2 = 1+1+0+1+1 = 4.000$;
$\operatorname{RSE}=\sqrt{\operatorname{RSS}/(n-2)}=\sqrt{4/3}=1.155$ (thousand EUR).
\emph{Interpretation:} even if the true line were known, observed weekly sales would deviate
from it by about $1.155$ thousand EUR ($\approx 1{,}155$ EUR) on average --- an estimate of
$\sigma$, the standard deviation of $\varepsilon$.

\medskip
\textbf{(d) Standard error, $t$-test and confidence interval.}
\[
\operatorname{SE}(\hat{\beta}_1)=\frac{1.1547}{\sqrt{10}}=\frac{1.1547}{3.1623}=0.365,\qquad
t=\frac{\hat{\beta}_1-0}{\operatorname{SE}(\hat{\beta}_1)}=\frac{3.000}{0.3651}=8.216
\quad(\text{df}=n-2=3).
\]
Since $|t|=8.216 > 3.182=t_{0.975,3}$, reject $H_0$ at the 5\,\% level.
95\,\% confidence interval:
\[
3.000 \pm 3.182\cdot 0.3651 = 3.000 \pm 1.162 = [1.838,\; 4.162].
\]
Zero is not contained in the interval --- consistent with the test.
\emph{Conclusion:} there is a statistically significant (positive) linear association between
the number of radio spots and weekly sales.

\medskip
\textbf{(e) TSS and $R^2$.}
$\operatorname{TSS}=\sum_i (y_i-\bar{y})^2 = 25+16+0+4+49 = 94.000$;
\[
R^2=1-\frac{\operatorname{RSS}}{\operatorname{TSS}}=1-\frac{4}{94}=\frac{90}{94}=0.957.
\]
About $95.7\,\%$ of the variance of weekly sales is explained by the linear regression on the
number of radio spots. (In simple regression $R^2=r^2$, so $r=\sqrt{0.957}\approx 0.979$.)

\medskip
\textbf{(f) Prediction at $x=6$.}
$\hat{y}(6)=2+3\cdot 6=20.000$ thousand EUR ($=20{,}000$ EUR).
Caution: $x=6$ lies \emph{outside} the observed range $x\in[1,5]$ --- an extrapolation; the
linear relationship has not been observed there and may not hold. (Also, $n=5$ is very small.)
\end{solutionbox}
\fi

% ============================================================
\problem{5: Multiple linear regression --- interpreting Python output}{24}

A consultancy analyses the salaries of $n=60$ employees of a software company.
The variables are \texttt{salary} (annual salary in 1{,}000 EUR), \texttt{experience}
(professional experience in years) and the dummy variable \texttt{female}
($=1$ for female employees, $=0$ for male employees). The model
\[
\texttt{salary}=\beta_0+\beta_1\,\texttt{experience}+\beta_2\,\texttt{female}
+\beta_3\,(\texttt{experience}\times\texttt{female})+\varepsilon
\]
was fitted by least squares in Python (\texttt{statsmodels}), with the following output:

\begin{center}
\begin{tabular}{lrrrr}
\toprule
 & coef & std err & $t$ & $P>|t|$\\
\midrule
Intercept                    & 46.2000 & 1.100 & 42.000 & $<$0.0001\\
experience                   &  2.3500 & 0.250 &  9.400 & $<$0.0001\\
female                       & $-3.1000$ & 1.550 & $-2.000$ & 0.0504\\
experience:female            & $-0.8500$ & 0.350 & $-2.429$ & 0.0184\\
\bottomrule
\end{tabular}

\medskip
$n=60$; \quad $R^2 = 0.712$; \quad adjusted $R^2 = 0.697$;\\[0.5mm]
$F$-statistic $= 46.15$ on 3 and 56 degrees of freedom, \ Prob($F$-statistic) $< 0.0001$.
\end{center}

\begin{enumerate}[label=(\alph*),leftmargin=7mm,itemsep=2.5mm]
  \item \pp{6} Interpret all four estimated coefficients precisely in the context of the
    problem (with units), paying particular attention to the dummy variable and the
    interaction term.
  \item \pp{4} Which coefficients are significantly different from zero at level
    $\alpha=0.05$? Verify for the interaction term that the reported $t$-value is consistent
    with the reported coefficient and standard error. Should \texttt{female} be removed from
    the model because of its $p$-value? Discuss briefly.
  \item \pp{4} Write down the fitted regression equation separately for men and for women.
  \item \pp{4} Predict the annual salary of (i) a man with 10 years of experience and (ii) a
    woman with 10 years of experience. What is the estimated salary gap between women and men
    at 10 years of experience, and how can it be written in terms of $\hat{\beta}_2$ and
    $\hat{\beta}_3$?
  \item \pp{3} Interpret $R^2$. State the null hypothesis of the reported $F$-test and the
    conclusion that follows from the output.
  \item \pp{3} Model diagnostics:
    \begin{enumerate}[label=(\roman*),nosep,topsep=1mm]
      \item A plot of the residuals against the fitted values shows a clear U-shape
        (positive residuals for small and large fitted values, negative in between).
        What does this indicate, and name one possible remedy.
      \item In an extended model, the predictor \texttt{age} has a variance inflation factor
        of $\operatorname{VIF}=12$. What does this mean, and what is the practical consequence
        for the regression results?
    \end{enumerate}
\end{enumerate}

\ifdefined\withsolutions
\begin{solutionbox}
\textbf{(a) Interpretation of the coefficients} (salary in 1{,}000 EUR per year):
\begin{itemize}[nosep,leftmargin=5mm]
  \item \emph{Intercept} $46.20$: expected annual salary of a \emph{man}
    (\texttt{female}=0) with 0 years of experience: $46{,}200$ EUR.
  \item \emph{experience} $2.35$: for \emph{men}, each additional year of experience increases
    the expected salary by $2{,}350$ EUR. (Because of the interaction, this slope refers to the
    baseline group \texttt{female}=0 only.)
  \item \emph{female} $-3.10$: at 0 years of experience, women are expected to earn
    $3{,}100$ EUR \emph{less} than men --- the difference between the two intercepts.
  \item \emph{experience:female} $-0.85$: the experience slope for women is $0.85$ \emph{lower}
    than for men: women gain $2.35-0.85=1.50$ thousand EUR per additional year. Equivalently,
    the expected salary gap between men and women \emph{widens} by $850$ EUR with every year
    of experience.
\end{itemize}

\medskip
\textbf{(b) Significance at $\alpha=0.05$.}
Intercept and \texttt{experience}: $p<0.0001$ $\Rightarrow$ highly significant.
\texttt{experience:female}: $p=0.0184<0.05$ $\Rightarrow$ significant.
Consistency check: $t=\hat{\beta}_3/\operatorname{SE}(\hat{\beta}_3)=-0.85/0.35=-2.429$
\checkmark\ (matches the output; $p$-value from $t_{56}$).
\texttt{female}: $p=0.0504>0.05$ $\Rightarrow$ \emph{not} significant at exactly the 5\,\%
level, though only marginally. It should nevertheless \emph{not} be removed: by the
\emph{hierarchy principle}, if an interaction ($\texttt{experience}\times\texttt{female}$) is
kept in the model, the corresponding main effects should be kept as well --- otherwise the
interaction changes its meaning. Moreover, $p=0.0504$ is borderline evidence, not evidence of
no effect.

\medskip
\textbf{(c) Fitted equations per group.}
\begin{align*}
\text{Men } (\texttt{female}=0):\quad & \widehat{\texttt{salary}} = 46.20 + 2.35\,\texttt{experience}\\
\text{Women } (\texttt{female}=1):\quad & \widehat{\texttt{salary}} = (46.20-3.10) + (2.35-0.85)\,\texttt{experience}\\
 & \phantom{\widehat{\texttt{salary}}} = 43.10 + 1.50\,\texttt{experience}
\end{align*}
Two different intercepts \emph{and} two different slopes (that is exactly what the interaction allows).

\medskip
\textbf{(d) Predictions at 10 years of experience.}
\begin{align*}
\text{Man: } & 46.20 + 2.35\cdot 10 = 69.70 \;\Rightarrow\; 69{,}700 \text{ EUR}\\
\text{Woman: } & 43.10 + 1.50\cdot 10 = 58.10 \;\Rightarrow\; 58{,}100 \text{ EUR}
\end{align*}
Estimated gap: $58.10-69.70=-11.60$, i.e.\ women earn $11{,}600$ EUR less. In terms of the
coefficients: $\hat{\beta}_2+10\,\hat{\beta}_3=-3.10-8.50=-11.60$ thousand EUR.

\medskip
\textbf{(e) $R^2$ and $F$-test.}
$R^2=0.712$: the three model terms explain $71.2\,\%$ of the variance of salaries in the sample.
$F$-test: $H_0\!: \beta_1=\beta_2=\beta_3=0$ (none of the predictors is related to salary)
against the alternative that at least one $\beta_j\neq 0$. With $F=46.15$ (3 and 56 df) and
$p<0.0001$, $H_0$ is clearly rejected: the model is useful overall.

\medskip
\textbf{(f) Diagnostics.}
(i) A U-shaped residual-vs-fitted plot indicates \emph{non-linearity}: the linear model misses
systematic curvature in the relationship. Remedy: include a non-linear term, e.g.\ add
$\texttt{experience}^2$ (polynomial regression), or transform variables
(e.g.\ $\log(\texttt{salary})$).
(ii) $\operatorname{VIF}_j = 1/(1-R^2_{X_j\mid X_{-j}})$; $\operatorname{VIF}=12$ means that
regressing \texttt{age} on the other predictors gives $R^2\approx 0.917$: \texttt{age} is
almost a linear combination of the other predictors (strong \emph{multicollinearity}; common
rules of thumb flag VIF $>5$ or $>10$). Consequence: $\operatorname{SE}(\hat{\beta}_{\texttt{age}})$
is inflated by a factor $\sqrt{12}\approx 3.5$ relative to the uncorrelated case --- the
coefficient is estimated very imprecisely, its sign may flip between samples, and its $t$-test
loses power, even though the model's overall fit and predictions may be unaffected.
\end{solutionbox}
\fi

\vspace{6mm}
\begin{center}
\rule{0.35\textwidth}{0.4pt}\\[1mm]
\textit{End of the examination --- good luck!}
\end{center}

\end{document}
"
% ============================================================
\documentclass[11pt,a4paper]{article}
\usepackage[utf8]{inputenc}
\usepackage[T1]{fontenc}
\usepackage[english]{babel}
\usepackage[margin=2.2cm]{geometry}
\usepackage{amsmath,amssymb}
\usepackage{booktabs,array}
\usepackage{enumitem}
\usepackage{xcolor}
\usepackage{tcolorbox}
\tcbuselibrary{skins,breakable}

\setlength{\parindent}{0pt}

% Teal solution box (only shown when \withsolutions is defined)
\newtcolorbox{solutionbox}{enhanced,breakable,colback=teal!5,colframe=teal!55!black,
  boxrule=0pt,leftrule=2.5pt,arc=2pt,fonttitle=\bfseries,title=Solution,
  left=2mm,right=2mm,top=1mm,bottom=1.5mm}

\newcommand{\problem}[2]{\par\vspace{7mm}\noindent{\large\textbf{Problem #1}}%
  \hfill{\large\textbf{(#2 points)}}\par\nopagebreak\vspace{0.6mm}%
  \noindent\rule{\textwidth}{0.6pt}\par\nopagebreak\vspace{2.5mm}}
\newcommand{\pp}[1]{\textbf{[#1~P]}\enspace}

\begin{document}

% ------------------------------------------------------------
% Header block
% ------------------------------------------------------------
\begin{center}
  {\Large\textbf{HSBI --- Bielefeld University of Applied Sciences and Arts}}\\[1.2mm]
  {\large Quantitative Research Methods}\\[1mm]
  {\normalsize Summer Semester 2026 \quad$\cdot$\quad Examiner: Prof.\ Dr.\ Christoph Weisser}\\[4.5mm]
  {\LARGE\textbf{Mock Examination 1 (Lectures 1--4)}}\\[1.2mm]
  {\normalsize ISLP Chapters 1--3: Statistical Learning $\cdot$ $K$-Nearest Neighbours $\cdot$ Linear Regression}%
  \ifdefined\withsolutions\\[2.2mm]{\large\textcolor{teal!55!black}{\textbf{--- Version with detailed solutions ---}}}\fi
\end{center}

\vspace{2.5mm}
\begin{center}
\begin{tabular}{@{}ll@{}}
  \textbf{Duration:} & \textbf{90 minutes}\\
  \textbf{Total credit:} & \textbf{90 points}\\
  \textbf{Allowed aids:} & non-programmable calculator; one handwritten A4 sheet of notes (both sides)\\
\end{tabular}
\end{center}

\vspace{3mm}
Name: \rule[-1mm]{0.38\textwidth}{0.4pt}\hfill Matriculation no.: \rule[-1mm]{0.24\textwidth}{0.4pt}

\vspace{4.5mm}
\begin{center}
\renewcommand{\arraystretch}{1.35}
\begin{tabular}{|l|c|c|c|c|c|c|}
\hline
\textbf{Problem} & \textbf{P1} & \textbf{P2} & \textbf{P3} & \textbf{P4} & \textbf{P5} & \textbf{Total}\\\hline
Maximum points & 15 & 15 & 12 & 24 & 24 & \textbf{90}\\\hline
Achieved points & \hspace{9mm} & \hspace{9mm} & \hspace{9mm} & \hspace{9mm} & \hspace{9mm} & \hspace{11mm}\\\hline
\end{tabular}
\end{center}

\vspace{2.5mm}
\textbf{Instructions}
\begin{itemize}[nosep,leftmargin=5.5mm]
  \item Justify all answers. Numerical results without a derivation or a brief reasoning receive no credit.
  \item Round final numerical results to 3 decimal places unless stated otherwise. Small deviations caused by carrying rounded intermediate results are accepted.
  \item The number of points per problem roughly equals the recommended working time in minutes.
\end{itemize}
\vspace{1mm}\noindent\rule{\textwidth}{0.6pt}

% ============================================================
\problem{1: Fundamentals of statistical learning}{15}

\begin{enumerate}[label=(\alph*),leftmargin=7mm,itemsep=2.5mm]
  \item \pp{4} Explain the difference between using a statistical learning method for
    \emph{prediction} and for \emph{inference}. Give one concrete example of a question
    for each of the two goals.
  \item \pp{4} Contrast \emph{parametric} and \emph{non-parametric} methods for estimating
    the unknown function $f$ in $Y=f(X)+\varepsilon$: describe the two approaches briefly
    and state one advantage of each. Which general trade-off between \emph{flexibility}
    and \emph{interpretability} lies behind this choice?
  \item \pp{3} Classify each of the following scenarios as \emph{regression},
    \emph{classification} or \emph{unsupervised learning}, with a one-sentence justification:
    \begin{enumerate}[label=(\roman*),nosep,topsep=1mm]
      \item A telecom provider has monthly usage, contract length and age for 800 customers
        and wants to predict whether each customer will cancel the contract next month.
      \item An economist records profit, number of employees and industry for 90 firms in
        order to understand which of these factors are associated with CEO salary, and how strongly.
      \item An online retailer has the shopping baskets of 5{,}000 customers and wants to
        discover groups of customers with similar purchasing behaviour.
    \end{enumerate}
  \item \pp{4} Explain why the training MSE of a method is in general \emph{smaller} than its
    test MSE. What is \emph{overfitting}, and how does it show up when training and test MSE
    are plotted against model flexibility?
\end{enumerate}

\ifdefined\withsolutions
\begin{solutionbox}
\textbf{(a) Prediction vs.\ inference.}
\emph{Prediction:} we want an accurate output $\hat{Y}=\hat{f}(X)$ for new inputs $X$; the exact
form of $\hat{f}$ is secondary, it may be treated as a black box (accuracy is what counts).
Example: predict whether an incoming e-mail is spam, or predict next month's demand for a product.
\emph{Inference:} we want to understand \emph{how} $Y$ changes with $X_1,\dots,X_p$: which
predictors matter, in which direction and how strongly. Here $\hat{f}$ must be interpretable.
Example: which advertising channels increase sales, and by how much per euro spent?

\medskip
\textbf{(b) Parametric vs.\ non-parametric.}
\emph{Parametric:} a two-step approach: (1) assume a functional form for $f$, e.g.\ the linear
model $f(X)=\beta_0+\beta_1X_1+\dots+\beta_pX_p$; (2) estimate the finite set of parameters
(e.g.\ by least squares). The problem of estimating $f$ is reduced to estimating a few parameters.
Advantage: needs relatively few observations, is simple to estimate and easy to interpret.
\emph{Non-parametric:} no explicit functional form is assumed; $\hat{f}$ is constructed to follow
the data directly (e.g.\ $K$-nearest neighbours, splines). Advantage: can capture a much wider
range of possible shapes of $f$, since no potentially wrong form is imposed.
The underlying trade-off: more flexible methods can reduce bias but need many observations, risk
overfitting and are harder to interpret; if inference is the goal, more restrictive
(interpretable) models are usually preferred --- this is the \emph{flexibility--interpretability}
trade-off.

\medskip
\textbf{(c) Scenarios.}
(i) \emph{Classification} (supervised): the response ``cancels: yes/no'' is qualitative;
the goal is prediction.
(ii) \emph{Regression} (supervised): the response CEO salary is quantitative;
the goal is inference (which factors are associated with salary).
(iii) \emph{Unsupervised learning}: there is no response variable; the goal is to find
clusters of similar customers.

\medskip
\textbf{(d) Training vs.\ test MSE and overfitting.}
The method is fitted by making the \emph{training} error small, so $\hat{f}$ adapts to the
training observations --- including their random noise. The training MSE therefore
systematically underestimates the error on fresh data, and for a given data set
training MSE $\le$ test MSE in general. As flexibility increases, the training MSE decreases
monotonically, while the test MSE first decreases and then increases (U-shape).
\emph{Overfitting} means the method follows patterns in the training noise that do not exist
in the population; a less flexible method would achieve a \emph{smaller test} MSE. It shows up
as a widening gap: very small training MSE combined with a rising test MSE.
\end{solutionbox}
\fi

% ============================================================
\problem{2: Bias--variance decomposition}{15}

For the prediction of a response at a fixed test point $x_0$, three methods of increasing
flexibility were analysed. A simulation study yielded the following values
(all in squared units of $Y$); the noise variance is $\operatorname{Var}(\varepsilon)=\sigma^2=1.00$:

\begin{center}
\begin{tabular}{lccc}
\toprule
 & \textbf{Method A} & \textbf{Method B} & \textbf{Method C}\\
 & (linear) & (moderately flexible) & (very flexible)\\
\midrule
$[\operatorname{Bias}(\hat{f}(x_0))]^2$ & 2.25 & 0.16 & 0.01\\
$\operatorname{Var}(\hat{f}(x_0))$      & 0.25 & 1.44 & 3.60\\
\bottomrule
\end{tabular}
\end{center}

\begin{enumerate}[label=(\alph*),leftmargin=7mm,itemsep=2.5mm]
  \item \pp{4} State the decomposition of the expected test MSE
    $E\!\left[(y_0-\hat{f}(x_0))^2\right]$ into its three components and explain in one
    sentence each what the three components measure.
  \item \pp{5} Compute the expected test MSE at $x_0$ for each of the three methods.
    Which method should be chosen, and why?
  \item \pp{3} What is the smallest expected test MSE at $x_0$ that \emph{any} method
    (however good) could achieve here? Justify your answer.
  \item \pp{3} Describe how squared bias, variance, training MSE and expected test MSE
    typically behave as the flexibility of a method increases, and explain how the numbers
    in the table illustrate the characteristic U-shape of the test MSE.
\end{enumerate}

\ifdefined\withsolutions
\begin{solutionbox}
\textbf{(a) Decomposition.}
\[
E\big[(y_0-\hat{f}(x_0))^2\big]
=\operatorname{Var}(\hat{f}(x_0))
+\big[\operatorname{Bias}(\hat{f}(x_0))\big]^2
+\operatorname{Var}(\varepsilon).
\]
\emph{Variance:} how much $\hat{f}(x_0)$ would change if the method were fitted on a different
training set. \emph{Squared bias:} the systematic error made by approximating the true $f$
by the (possibly too simple) model. \emph{Irreducible error} $\sigma^2$: the noise variance
of $\varepsilon$, which no method can remove.

\medskip
\textbf{(b) Expected test MSE.}
\begin{align*}
\text{A: } & 2.25+0.25+1.00 = 3.50\\
\text{B: } & 0.16+1.44+1.00 = \mathbf{2.60}\\
\text{C: } & 0.01+3.60+1.00 = 4.61
\end{align*}
Choose \textbf{Method B}: it has the smallest expected test MSE. Its moderate flexibility buys a
large bias reduction (2.25 $\to$ 0.16) at an acceptable variance cost, whereas the very flexible
Method C pays more in variance ($+2.16$) than it gains in bias ($-0.15$).

\medskip
\textbf{(c) Lower bound.}
$\sigma^2=1.00$. Squared bias and variance are non-negative, so
$E[(y_0-\hat{f}(x_0))^2]\ge \operatorname{Var}(\varepsilon)=1.00$ for every method.
Even the true $f$ (zero bias, zero variance) cannot predict the noise term $\varepsilon_0$:
$1.00$ is the \emph{irreducible error}.

\medskip
\textbf{(d) Behaviour with increasing flexibility.}
Squared bias decreases (typically monotonically), variance increases, and the training MSE
decreases monotonically (a more flexible fit follows the training data ever more closely).
The expected test MSE is the sum bias$^2$ + variance + $\sigma^2$: it first falls, while the
bias reduction dominates, and then rises, when the variance increase dominates --- the U-shape.
In the table: $3.50 \to 2.60 \to 4.61$ for A $\to$ B $\to$ C; the minimum is attained at the
intermediate flexibility (B).
\end{solutionbox}
\fi

% ============================================================
\problem{3: $K$-nearest neighbours classification}{12}

An online shop wants to predict whether a visitor makes a purchase. For seven past visitors,
two standardised activity scores $X_1$ (pages viewed) and $X_2$ (minutes on site) and the
outcome $Y$ (purchase: Yes/No) were recorded:

\begin{center}
\begin{tabular}{cccc}
\toprule
Obs. & $X_1$ & $X_2$ & $Y$\\
\midrule
1 & 4 & 3 & Yes\\
2 & 2 & 2 & No\\
3 & 6 & 2 & No\\
4 & 4 & 5 & No\\
5 & 1 & 1 & Yes\\
6 & 7 & 4 & Yes\\
7 & 5 & 6 & No\\
\bottomrule
\end{tabular}
\end{center}

A new visitor has scores $x_0=(X_1,X_2)=(4,2)$.

\begin{enumerate}[label=(\alph*),leftmargin=7mm,itemsep=2.5mm]
  \item \pp{4} Compute the Euclidean distance between $x_0$ and each of the seven
    observations (3 decimals; you may rank by squared distances). How is the new visitor
    classified with $K=1$?
  \item \pp{4} How is the new visitor classified with $K=3$? State also the estimated
    conditional probability $\widehat{\Pr}(Y=\text{Yes}\mid X=x_0)$ that the $K=3$
    classifier uses.
  \item \pp{2} How does the choice of $K$ affect the bias and the variance of the KNN
    classifier (small $K$ vs.\ large $K$)? What is the \emph{training} error rate of the
    1-nearest-neighbour classifier?
  \item \pp{2} What is the \emph{Bayes classifier}, and why can it not be used directly in
    practice? In what sense is KNN related to it?
\end{enumerate}

\ifdefined\withsolutions
\begin{solutionbox}
\textbf{(a) Distances and $K=1$.} With $d(x_0,x_i)=\sqrt{(X_{1i}-4)^2+(X_{2i}-2)^2}$:
\begin{center}
\begin{tabular}{lccccccc}
\toprule
Obs. & 1 & 2 & 3 & 4 & 5 & 6 & 7\\
\midrule
$d^2$ & 1 & 4 & 4 & 9 & 10 & 13 & 17\\
$d$ & 1.000 & 2.000 & 2.000 & 3.000 & 3.162 & 3.606 & 4.123\\
$Y$ & Yes & No & No & No & Yes & Yes & No\\
\bottomrule
\end{tabular}
\end{center}
$K=1$: the single nearest neighbour is Obs.~1 ($d=1.000$) with label Yes
$\Rightarrow$ classify as \textbf{Yes}.

\medskip
\textbf{(b) $K=3$.} The three nearest neighbours are Obs.~1 ($d=1.000$, Yes),
Obs.~2 ($d=2.000$, No) and Obs.~3 ($d=2.000$, No). Estimated conditional probabilities:
$\widehat{\Pr}(\text{Yes}\mid x_0)=1/3\approx 0.333$ and
$\widehat{\Pr}(\text{No}\mid x_0)=2/3\approx 0.667$.
Majority vote $\Rightarrow$ classify as \textbf{No}. (Note that $K=1$ and $K=3$ disagree ---
the decision depends on the flexibility of the classifier.)

\medskip
\textbf{(c) Effect of $K$.} Small $K$ (e.g.\ $K=1$): very flexible, jagged decision boundary
$\Rightarrow$ \emph{low bias, high variance}. Large $K$: averaging over many neighbours gives a
smooth, almost linear boundary $\Rightarrow$ \emph{high bias, low variance}. The training error
rate of 1-NN is $0$ (ignoring ties): each training point is its own nearest neighbour --- which
says nothing about its test error.

\medskip
\textbf{(d) Bayes classifier.} It assigns each $x$ to the class $j$ with the largest \emph{true}
conditional probability $\Pr(Y=j\mid X=x)$; it minimises the expected test error rate, and its
error (the \emph{Bayes error rate}) is the lower bound for all classifiers. It is unattainable
in practice because the true conditional distribution of $Y$ given $X$ is unknown. KNN mimics
it by \emph{estimating} these conditional probabilities locally from the $K$ nearest training
observations and then classifying to the largest estimated probability.
\end{solutionbox}
\fi

% ============================================================
\problem{4: Simple linear regression by hand}{24}

A caf\'e chain investigates the effect of local radio advertising on sales.
For $n=5$ weeks it records the number of radio spots per week $x$ and the weekly sales $y$
(in 1{,}000 EUR):

\begin{center}
\begin{tabular}{lccccc}
\toprule
Week $i$ & 1 & 2 & 3 & 4 & 5\\
\midrule
Spots $x_i$ & 1 & 2 & 3 & 4 & 5\\
Sales $y_i$ & 6 & 7 & 11 & 13 & 18\\
\bottomrule
\end{tabular}
\end{center}

You may use without proof: $\sum x_i = 15$, $\sum y_i = 55$, $\sum x_i^2 = 55$,
$\sum x_i y_i = 195$. Consider the model $Y=\beta_0+\beta_1 X+\varepsilon$.

\begin{enumerate}[label=(\alph*),leftmargin=7mm,itemsep=2.5mm]
  \item \pp{6} Compute the least-squares estimates $\hat{\beta}_1$ and $\hat{\beta}_0$, and
    write down the fitted regression line. Interpret both estimates in the context of the
    problem (with units).
  \item \pp{4} Compute the fitted values $\hat{y}_i$ and the residuals $e_i$ for all five
    weeks, and verify that $\sum_{i=1}^{5} e_i = 0$.
  \item \pp{3} Compute the residual sum of squares (RSS) and the residual standard error
    (RSE), and interpret the RSE.
  \item \pp{6} Compute $\operatorname{SE}(\hat{\beta}_1)=\operatorname{RSE}/\sqrt{\sum_i (x_i-\bar{x})^2}$
    and test $H_0\!:\beta_1=0$ against $H_1\!:\beta_1\neq 0$ at level $\alpha=0.05$ using the
    $t$-statistic. Also give the 95\,\% confidence interval for $\beta_1$.
    \emph{Use $t_{0.975,\,3}=3.182$ (97.5\,\% quantile of the $t$-distribution with 3 degrees
    of freedom).} State your conclusion in one sentence.
  \item \pp{3} Compute the total sum of squares (TSS) and the coefficient of determination
    $R^2$, and interpret $R^2$.
  \item \pp{2} Predict the weekly sales for a week with $x=6$ radio spots. Name one reason
    why this particular prediction should be treated with caution.
\end{enumerate}

\ifdefined\withsolutions
\begin{solutionbox}
\textbf{(a) Coefficient estimates.} $\bar{x}=15/5=3$, $\bar{y}=55/5=11$.
\[
S_{xy}=\sum x_iy_i-n\bar{x}\bar{y}=195-5\cdot 3\cdot 11=30,\qquad
S_{xx}=\sum x_i^2-n\bar{x}^2=55-5\cdot 9=10.
\]
\[
\hat{\beta}_1=\frac{S_{xy}}{S_{xx}}=\frac{30}{10}=3.000,\qquad
\hat{\beta}_0=\bar{y}-\hat{\beta}_1\bar{x}=11-3\cdot 3=2.000.
\]
Fitted line: $\hat{y}=2.000+3.000\,x$.
\emph{Interpretation:} each additional radio spot per week is associated with an expected
increase in weekly sales of $3.000$ thousand EUR ($=3{,}000$ EUR). The intercept $2.000$ is the
expected weekly sales ($2{,}000$ EUR) in a week without any spots --- a mild extrapolation,
since $x=0$ was not observed ($x$ ranges from 1 to 5).

\medskip
\textbf{(b) Fitted values and residuals.} $\hat{y}_i=2+3x_i$ and $e_i=y_i-\hat{y}_i$:
\begin{center}
\begin{tabular}{lccccc}
\toprule
$i$ & 1 & 2 & 3 & 4 & 5\\
\midrule
$\hat{y}_i$ & 5 & 8 & 11 & 14 & 17\\
$e_i$ & $1$ & $-1$ & $0$ & $-1$ & $1$\\
\bottomrule
\end{tabular}
\end{center}
$\sum_i e_i = 1-1+0-1+1 = 0$ \checkmark\ (this always holds for least squares with an intercept).

\medskip
\textbf{(c) RSS and RSE.}
$\operatorname{RSS}=\sum_i e_i^2 = 1+1+0+1+1 = 4.000$;
$\operatorname{RSE}=\sqrt{\operatorname{RSS}/(n-2)}=\sqrt{4/3}=1.155$ (thousand EUR).
\emph{Interpretation:} even if the true line were known, observed weekly sales would deviate
from it by about $1.155$ thousand EUR ($\approx 1{,}155$ EUR) on average --- an estimate of
$\sigma$, the standard deviation of $\varepsilon$.

\medskip
\textbf{(d) Standard error, $t$-test and confidence interval.}
\[
\operatorname{SE}(\hat{\beta}_1)=\frac{1.1547}{\sqrt{10}}=\frac{1.1547}{3.1623}=0.365,\qquad
t=\frac{\hat{\beta}_1-0}{\operatorname{SE}(\hat{\beta}_1)}=\frac{3.000}{0.3651}=8.216
\quad(\text{df}=n-2=3).
\]
Since $|t|=8.216 > 3.182=t_{0.975,3}$, reject $H_0$ at the 5\,\% level.
95\,\% confidence interval:
\[
3.000 \pm 3.182\cdot 0.3651 = 3.000 \pm 1.162 = [1.838,\; 4.162].
\]
Zero is not contained in the interval --- consistent with the test.
\emph{Conclusion:} there is a statistically significant (positive) linear association between
the number of radio spots and weekly sales.

\medskip
\textbf{(e) TSS and $R^2$.}
$\operatorname{TSS}=\sum_i (y_i-\bar{y})^2 = 25+16+0+4+49 = 94.000$;
\[
R^2=1-\frac{\operatorname{RSS}}{\operatorname{TSS}}=1-\frac{4}{94}=\frac{90}{94}=0.957.
\]
About $95.7\,\%$ of the variance of weekly sales is explained by the linear regression on the
number of radio spots. (In simple regression $R^2=r^2$, so $r=\sqrt{0.957}\approx 0.979$.)

\medskip
\textbf{(f) Prediction at $x=6$.}
$\hat{y}(6)=2+3\cdot 6=20.000$ thousand EUR ($=20{,}000$ EUR).
Caution: $x=6$ lies \emph{outside} the observed range $x\in[1,5]$ --- an extrapolation; the
linear relationship has not been observed there and may not hold. (Also, $n=5$ is very small.)
\end{solutionbox}
\fi

% ============================================================
\problem{5: Multiple linear regression --- interpreting Python output}{24}

A consultancy analyses the salaries of $n=60$ employees of a software company.
The variables are \texttt{salary} (annual salary in 1{,}000 EUR), \texttt{experience}
(professional experience in years) and the dummy variable \texttt{female}
($=1$ for female employees, $=0$ for male employees). The model
\[
\texttt{salary}=\beta_0+\beta_1\,\texttt{experience}+\beta_2\,\texttt{female}
+\beta_3\,(\texttt{experience}\times\texttt{female})+\varepsilon
\]
was fitted by least squares in Python (\texttt{statsmodels}), with the following output:

\begin{center}
\begin{tabular}{lrrrr}
\toprule
 & coef & std err & $t$ & $P>|t|$\\
\midrule
Intercept                    & 46.2000 & 1.100 & 42.000 & $<$0.0001\\
experience                   &  2.3500 & 0.250 &  9.400 & $<$0.0001\\
female                       & $-3.1000$ & 1.550 & $-2.000$ & 0.0504\\
experience:female            & $-0.8500$ & 0.350 & $-2.429$ & 0.0184\\
\bottomrule
\end{tabular}

\medskip
$n=60$; \quad $R^2 = 0.712$; \quad adjusted $R^2 = 0.697$;\\[0.5mm]
$F$-statistic $= 46.15$ on 3 and 56 degrees of freedom, \ Prob($F$-statistic) $< 0.0001$.
\end{center}

\begin{enumerate}[label=(\alph*),leftmargin=7mm,itemsep=2.5mm]
  \item \pp{6} Interpret all four estimated coefficients precisely in the context of the
    problem (with units), paying particular attention to the dummy variable and the
    interaction term.
  \item \pp{4} Which coefficients are significantly different from zero at level
    $\alpha=0.05$? Verify for the interaction term that the reported $t$-value is consistent
    with the reported coefficient and standard error. Should \texttt{female} be removed from
    the model because of its $p$-value? Discuss briefly.
  \item \pp{4} Write down the fitted regression equation separately for men and for women.
  \item \pp{4} Predict the annual salary of (i) a man with 10 years of experience and (ii) a
    woman with 10 years of experience. What is the estimated salary gap between women and men
    at 10 years of experience, and how can it be written in terms of $\hat{\beta}_2$ and
    $\hat{\beta}_3$?
  \item \pp{3} Interpret $R^2$. State the null hypothesis of the reported $F$-test and the
    conclusion that follows from the output.
  \item \pp{3} Model diagnostics:
    \begin{enumerate}[label=(\roman*),nosep,topsep=1mm]
      \item A plot of the residuals against the fitted values shows a clear U-shape
        (positive residuals for small and large fitted values, negative in between).
        What does this indicate, and name one possible remedy.
      \item In an extended model, the predictor \texttt{age} has a variance inflation factor
        of $\operatorname{VIF}=12$. What does this mean, and what is the practical consequence
        for the regression results?
    \end{enumerate}
\end{enumerate}

\ifdefined\withsolutions
\begin{solutionbox}
\textbf{(a) Interpretation of the coefficients} (salary in 1{,}000 EUR per year):
\begin{itemize}[nosep,leftmargin=5mm]
  \item \emph{Intercept} $46.20$: expected annual salary of a \emph{man}
    (\texttt{female}=0) with 0 years of experience: $46{,}200$ EUR.
  \item \emph{experience} $2.35$: for \emph{men}, each additional year of experience increases
    the expected salary by $2{,}350$ EUR. (Because of the interaction, this slope refers to the
    baseline group \texttt{female}=0 only.)
  \item \emph{female} $-3.10$: at 0 years of experience, women are expected to earn
    $3{,}100$ EUR \emph{less} than men --- the difference between the two intercepts.
  \item \emph{experience:female} $-0.85$: the experience slope for women is $0.85$ \emph{lower}
    than for men: women gain $2.35-0.85=1.50$ thousand EUR per additional year. Equivalently,
    the expected salary gap between men and women \emph{widens} by $850$ EUR with every year
    of experience.
\end{itemize}

\medskip
\textbf{(b) Significance at $\alpha=0.05$.}
Intercept and \texttt{experience}: $p<0.0001$ $\Rightarrow$ highly significant.
\texttt{experience:female}: $p=0.0184<0.05$ $\Rightarrow$ significant.
Consistency check: $t=\hat{\beta}_3/\operatorname{SE}(\hat{\beta}_3)=-0.85/0.35=-2.429$
\checkmark\ (matches the output; $p$-value from $t_{56}$).
\texttt{female}: $p=0.0504>0.05$ $\Rightarrow$ \emph{not} significant at exactly the 5\,\%
level, though only marginally. It should nevertheless \emph{not} be removed: by the
\emph{hierarchy principle}, if an interaction ($\texttt{experience}\times\texttt{female}$) is
kept in the model, the corresponding main effects should be kept as well --- otherwise the
interaction changes its meaning. Moreover, $p=0.0504$ is borderline evidence, not evidence of
no effect.

\medskip
\textbf{(c) Fitted equations per group.}
\begin{align*}
\text{Men } (\texttt{female}=0):\quad & \widehat{\texttt{salary}} = 46.20 + 2.35\,\texttt{experience}\\
\text{Women } (\texttt{female}=1):\quad & \widehat{\texttt{salary}} = (46.20-3.10) + (2.35-0.85)\,\texttt{experience}\\
 & \phantom{\widehat{\texttt{salary}}} = 43.10 + 1.50\,\texttt{experience}
\end{align*}
Two different intercepts \emph{and} two different slopes (that is exactly what the interaction allows).

\medskip
\textbf{(d) Predictions at 10 years of experience.}
\begin{align*}
\text{Man: } & 46.20 + 2.35\cdot 10 = 69.70 \;\Rightarrow\; 69{,}700 \text{ EUR}\\
\text{Woman: } & 43.10 + 1.50\cdot 10 = 58.10 \;\Rightarrow\; 58{,}100 \text{ EUR}
\end{align*}
Estimated gap: $58.10-69.70=-11.60$, i.e.\ women earn $11{,}600$ EUR less. In terms of the
coefficients: $\hat{\beta}_2+10\,\hat{\beta}_3=-3.10-8.50=-11.60$ thousand EUR.

\medskip
\textbf{(e) $R^2$ and $F$-test.}
$R^2=0.712$: the three model terms explain $71.2\,\%$ of the variance of salaries in the sample.
$F$-test: $H_0\!: \beta_1=\beta_2=\beta_3=0$ (none of the predictors is related to salary)
against the alternative that at least one $\beta_j\neq 0$. With $F=46.15$ (3 and 56 df) and
$p<0.0001$, $H_0$ is clearly rejected: the model is useful overall.

\medskip
\textbf{(f) Diagnostics.}
(i) A U-shaped residual-vs-fitted plot indicates \emph{non-linearity}: the linear model misses
systematic curvature in the relationship. Remedy: include a non-linear term, e.g.\ add
$\texttt{experience}^2$ (polynomial regression), or transform variables
(e.g.\ $\log(\texttt{salary})$).
(ii) $\operatorname{VIF}_j = 1/(1-R^2_{X_j\mid X_{-j}})$; $\operatorname{VIF}=12$ means that
regressing \texttt{age} on the other predictors gives $R^2\approx 0.917$: \texttt{age} is
almost a linear combination of the other predictors (strong \emph{multicollinearity}; common
rules of thumb flag VIF $>5$ or $>10$). Consequence: $\operatorname{SE}(\hat{\beta}_{\texttt{age}})$
is inflated by a factor $\sqrt{12}\approx 3.5$ relative to the uncorrelated case --- the
coefficient is estimated very imprecisely, its sign may flip between samples, and its $t$-test
loses power, even though the model's overall fit and predictions may be unaffected.
\end{solutionbox}
\fi

\vspace{6mm}
\begin{center}
\rule{0.35\textwidth}{0.4pt}\\[1mm]
\textit{End of the examination --- good luck!}
\end{center}

\end{document}
"
% ============================================================
\documentclass[11pt,a4paper]{article}
\usepackage[utf8]{inputenc}
\usepackage[T1]{fontenc}
\usepackage[english]{babel}
\usepackage[margin=2.2cm]{geometry}
\usepackage{amsmath,amssymb}
\usepackage{booktabs,array}
\usepackage{enumitem}
\usepackage{xcolor}
\usepackage{tcolorbox}
\tcbuselibrary{skins,breakable}

\setlength{\parindent}{0pt}

% Teal solution box (only shown when \withsolutions is defined)
\newtcolorbox{solutionbox}{enhanced,breakable,colback=teal!5,colframe=teal!55!black,
  boxrule=0pt,leftrule=2.5pt,arc=2pt,fonttitle=\bfseries,title=Solution,
  left=2mm,right=2mm,top=1mm,bottom=1.5mm}

\newcommand{\problem}[2]{\par\vspace{7mm}\noindent{\large\textbf{Problem #1}}%
  \hfill{\large\textbf{(#2 points)}}\par\nopagebreak\vspace{0.6mm}%
  \noindent\rule{\textwidth}{0.6pt}\par\nopagebreak\vspace{2.5mm}}
\newcommand{\pp}[1]{\textbf{[#1~P]}\enspace}

\begin{document}

% ------------------------------------------------------------
% Header block
% ------------------------------------------------------------
\begin{center}
  {\Large\textbf{HSBI --- Bielefeld University of Applied Sciences and Arts}}\\[1.2mm]
  {\large Quantitative Research Methods}\\[1mm]
  {\normalsize Summer Semester 2026 \quad$\cdot$\quad Examiner: Prof.\ Dr.\ Christoph Weisser}\\[4.5mm]
  {\LARGE\textbf{Mock Examination 1 (Lectures 1--4)}}\\[1.2mm]
  {\normalsize ISLP Chapters 1--3: Statistical Learning $\cdot$ $K$-Nearest Neighbours $\cdot$ Linear Regression}%
  \ifdefined\withsolutions\\[2.2mm]{\large\textcolor{teal!55!black}{\textbf{--- Version with detailed solutions ---}}}\fi
\end{center}

\vspace{2.5mm}
\begin{center}
\begin{tabular}{@{}ll@{}}
  \textbf{Duration:} & \textbf{90 minutes}\\
  \textbf{Total credit:} & \textbf{90 points}\\
  \textbf{Allowed aids:} & non-programmable calculator; one handwritten A4 sheet of notes (both sides)\\
\end{tabular}
\end{center}

\vspace{3mm}
Name: \rule[-1mm]{0.38\textwidth}{0.4pt}\hfill Matriculation no.: \rule[-1mm]{0.24\textwidth}{0.4pt}

\vspace{4.5mm}
\begin{center}
\renewcommand{\arraystretch}{1.35}
\begin{tabular}{|l|c|c|c|c|c|c|}
\hline
\textbf{Problem} & \textbf{P1} & \textbf{P2} & \textbf{P3} & \textbf{P4} & \textbf{P5} & \textbf{Total}\\\hline
Maximum points & 15 & 15 & 12 & 24 & 24 & \textbf{90}\\\hline
Achieved points & \hspace{9mm} & \hspace{9mm} & \hspace{9mm} & \hspace{9mm} & \hspace{9mm} & \hspace{11mm}\\\hline
\end{tabular}
\end{center}

\vspace{2.5mm}
\textbf{Instructions}
\begin{itemize}[nosep,leftmargin=5.5mm]
  \item Justify all answers. Numerical results without a derivation or a brief reasoning receive no credit.
  \item Round final numerical results to 3 decimal places unless stated otherwise. Small deviations caused by carrying rounded intermediate results are accepted.
  \item The number of points per problem roughly equals the recommended working time in minutes.
\end{itemize}
\vspace{1mm}\noindent\rule{\textwidth}{0.6pt}

% ============================================================
\problem{1: Fundamentals of statistical learning}{15}

\begin{enumerate}[label=(\alph*),leftmargin=7mm,itemsep=2.5mm]
  \item \pp{4} Explain the difference between using a statistical learning method for
    \emph{prediction} and for \emph{inference}. Give one concrete example of a question
    for each of the two goals.
  \item \pp{4} Contrast \emph{parametric} and \emph{non-parametric} methods for estimating
    the unknown function $f$ in $Y=f(X)+\varepsilon$: describe the two approaches briefly
    and state one advantage of each. Which general trade-off between \emph{flexibility}
    and \emph{interpretability} lies behind this choice?
  \item \pp{3} Classify each of the following scenarios as \emph{regression},
    \emph{classification} or \emph{unsupervised learning}, with a one-sentence justification:
    \begin{enumerate}[label=(\roman*),nosep,topsep=1mm]
      \item A telecom provider has monthly usage, contract length and age for 800 customers
        and wants to predict whether each customer will cancel the contract next month.
      \item An economist records profit, number of employees and industry for 90 firms in
        order to understand which of these factors are associated with CEO salary, and how strongly.
      \item An online retailer has the shopping baskets of 5{,}000 customers and wants to
        discover groups of customers with similar purchasing behaviour.
    \end{enumerate}
  \item \pp{4} Explain why the training MSE of a method is in general \emph{smaller} than its
    test MSE. What is \emph{overfitting}, and how does it show up when training and test MSE
    are plotted against model flexibility?
\end{enumerate}

\ifdefined\withsolutions
\begin{solutionbox}
\textbf{(a) Prediction vs.\ inference.}
\emph{Prediction:} we want an accurate output $\hat{Y}=\hat{f}(X)$ for new inputs $X$; the exact
form of $\hat{f}$ is secondary, it may be treated as a black box (accuracy is what counts).
Example: predict whether an incoming e-mail is spam, or predict next month's demand for a product.
\emph{Inference:} we want to understand \emph{how} $Y$ changes with $X_1,\dots,X_p$: which
predictors matter, in which direction and how strongly. Here $\hat{f}$ must be interpretable.
Example: which advertising channels increase sales, and by how much per euro spent?

\medskip
\textbf{(b) Parametric vs.\ non-parametric.}
\emph{Parametric:} a two-step approach: (1) assume a functional form for $f$, e.g.\ the linear
model $f(X)=\beta_0+\beta_1X_1+\dots+\beta_pX_p$; (2) estimate the finite set of parameters
(e.g.\ by least squares). The problem of estimating $f$ is reduced to estimating a few parameters.
Advantage: needs relatively few observations, is simple to estimate and easy to interpret.
\emph{Non-parametric:} no explicit functional form is assumed; $\hat{f}$ is constructed to follow
the data directly (e.g.\ $K$-nearest neighbours, splines). Advantage: can capture a much wider
range of possible shapes of $f$, since no potentially wrong form is imposed.
The underlying trade-off: more flexible methods can reduce bias but need many observations, risk
overfitting and are harder to interpret; if inference is the goal, more restrictive
(interpretable) models are usually preferred --- this is the \emph{flexibility--interpretability}
trade-off.

\medskip
\textbf{(c) Scenarios.}
(i) \emph{Classification} (supervised): the response ``cancels: yes/no'' is qualitative;
the goal is prediction.
(ii) \emph{Regression} (supervised): the response CEO salary is quantitative;
the goal is inference (which factors are associated with salary).
(iii) \emph{Unsupervised learning}: there is no response variable; the goal is to find
clusters of similar customers.

\medskip
\textbf{(d) Training vs.\ test MSE and overfitting.}
The method is fitted by making the \emph{training} error small, so $\hat{f}$ adapts to the
training observations --- including their random noise. The training MSE therefore
systematically underestimates the error on fresh data, and for a given data set
training MSE $\le$ test MSE in general. As flexibility increases, the training MSE decreases
monotonically, while the test MSE first decreases and then increases (U-shape).
\emph{Overfitting} means the method follows patterns in the training noise that do not exist
in the population; a less flexible method would achieve a \emph{smaller test} MSE. It shows up
as a widening gap: very small training MSE combined with a rising test MSE.
\end{solutionbox}
\fi

% ============================================================
\problem{2: Bias--variance decomposition}{15}

For the prediction of a response at a fixed test point $x_0$, three methods of increasing
flexibility were analysed. A simulation study yielded the following values
(all in squared units of $Y$); the noise variance is $\operatorname{Var}(\varepsilon)=\sigma^2=1.00$:

\begin{center}
\begin{tabular}{lccc}
\toprule
 & \textbf{Method A} & \textbf{Method B} & \textbf{Method C}\\
 & (linear) & (moderately flexible) & (very flexible)\\
\midrule
$[\operatorname{Bias}(\hat{f}(x_0))]^2$ & 2.25 & 0.16 & 0.01\\
$\operatorname{Var}(\hat{f}(x_0))$      & 0.25 & 1.44 & 3.60\\
\bottomrule
\end{tabular}
\end{center}

\begin{enumerate}[label=(\alph*),leftmargin=7mm,itemsep=2.5mm]
  \item \pp{4} State the decomposition of the expected test MSE
    $E\!\left[(y_0-\hat{f}(x_0))^2\right]$ into its three components and explain in one
    sentence each what the three components measure.
  \item \pp{5} Compute the expected test MSE at $x_0$ for each of the three methods.
    Which method should be chosen, and why?
  \item \pp{3} What is the smallest expected test MSE at $x_0$ that \emph{any} method
    (however good) could achieve here? Justify your answer.
  \item \pp{3} Describe how squared bias, variance, training MSE and expected test MSE
    typically behave as the flexibility of a method increases, and explain how the numbers
    in the table illustrate the characteristic U-shape of the test MSE.
\end{enumerate}

\ifdefined\withsolutions
\begin{solutionbox}
\textbf{(a) Decomposition.}
\[
E\big[(y_0-\hat{f}(x_0))^2\big]
=\operatorname{Var}(\hat{f}(x_0))
+\big[\operatorname{Bias}(\hat{f}(x_0))\big]^2
+\operatorname{Var}(\varepsilon).
\]
\emph{Variance:} how much $\hat{f}(x_0)$ would change if the method were fitted on a different
training set. \emph{Squared bias:} the systematic error made by approximating the true $f$
by the (possibly too simple) model. \emph{Irreducible error} $\sigma^2$: the noise variance
of $\varepsilon$, which no method can remove.

\medskip
\textbf{(b) Expected test MSE.}
\begin{align*}
\text{A: } & 2.25+0.25+1.00 = 3.50\\
\text{B: } & 0.16+1.44+1.00 = \mathbf{2.60}\\
\text{C: } & 0.01+3.60+1.00 = 4.61
\end{align*}
Choose \textbf{Method B}: it has the smallest expected test MSE. Its moderate flexibility buys a
large bias reduction (2.25 $\to$ 0.16) at an acceptable variance cost, whereas the very flexible
Method C pays more in variance ($+2.16$) than it gains in bias ($-0.15$).

\medskip
\textbf{(c) Lower bound.}
$\sigma^2=1.00$. Squared bias and variance are non-negative, so
$E[(y_0-\hat{f}(x_0))^2]\ge \operatorname{Var}(\varepsilon)=1.00$ for every method.
Even the true $f$ (zero bias, zero variance) cannot predict the noise term $\varepsilon_0$:
$1.00$ is the \emph{irreducible error}.

\medskip
\textbf{(d) Behaviour with increasing flexibility.}
Squared bias decreases (typically monotonically), variance increases, and the training MSE
decreases monotonically (a more flexible fit follows the training data ever more closely).
The expected test MSE is the sum bias$^2$ + variance + $\sigma^2$: it first falls, while the
bias reduction dominates, and then rises, when the variance increase dominates --- the U-shape.
In the table: $3.50 \to 2.60 \to 4.61$ for A $\to$ B $\to$ C; the minimum is attained at the
intermediate flexibility (B).
\end{solutionbox}
\fi

% ============================================================
\problem{3: $K$-nearest neighbours classification}{12}

An online shop wants to predict whether a visitor makes a purchase. For seven past visitors,
two standardised activity scores $X_1$ (pages viewed) and $X_2$ (minutes on site) and the
outcome $Y$ (purchase: Yes/No) were recorded:

\begin{center}
\begin{tabular}{cccc}
\toprule
Obs. & $X_1$ & $X_2$ & $Y$\\
\midrule
1 & 4 & 3 & Yes\\
2 & 2 & 2 & No\\
3 & 6 & 2 & No\\
4 & 4 & 5 & No\\
5 & 1 & 1 & Yes\\
6 & 7 & 4 & Yes\\
7 & 5 & 6 & No\\
\bottomrule
\end{tabular}
\end{center}

A new visitor has scores $x_0=(X_1,X_2)=(4,2)$.

\begin{enumerate}[label=(\alph*),leftmargin=7mm,itemsep=2.5mm]
  \item \pp{4} Compute the Euclidean distance between $x_0$ and each of the seven
    observations (3 decimals; you may rank by squared distances). How is the new visitor
    classified with $K=1$?
  \item \pp{4} How is the new visitor classified with $K=3$? State also the estimated
    conditional probability $\widehat{\Pr}(Y=\text{Yes}\mid X=x_0)$ that the $K=3$
    classifier uses.
  \item \pp{2} How does the choice of $K$ affect the bias and the variance of the KNN
    classifier (small $K$ vs.\ large $K$)? What is the \emph{training} error rate of the
    1-nearest-neighbour classifier?
  \item \pp{2} What is the \emph{Bayes classifier}, and why can it not be used directly in
    practice? In what sense is KNN related to it?
\end{enumerate}

\ifdefined\withsolutions
\begin{solutionbox}
\textbf{(a) Distances and $K=1$.} With $d(x_0,x_i)=\sqrt{(X_{1i}-4)^2+(X_{2i}-2)^2}$:
\begin{center}
\begin{tabular}{lccccccc}
\toprule
Obs. & 1 & 2 & 3 & 4 & 5 & 6 & 7\\
\midrule
$d^2$ & 1 & 4 & 4 & 9 & 10 & 13 & 17\\
$d$ & 1.000 & 2.000 & 2.000 & 3.000 & 3.162 & 3.606 & 4.123\\
$Y$ & Yes & No & No & No & Yes & Yes & No\\
\bottomrule
\end{tabular}
\end{center}
$K=1$: the single nearest neighbour is Obs.~1 ($d=1.000$) with label Yes
$\Rightarrow$ classify as \textbf{Yes}.

\medskip
\textbf{(b) $K=3$.} The three nearest neighbours are Obs.~1 ($d=1.000$, Yes),
Obs.~2 ($d=2.000$, No) and Obs.~3 ($d=2.000$, No). Estimated conditional probabilities:
$\widehat{\Pr}(\text{Yes}\mid x_0)=1/3\approx 0.333$ and
$\widehat{\Pr}(\text{No}\mid x_0)=2/3\approx 0.667$.
Majority vote $\Rightarrow$ classify as \textbf{No}. (Note that $K=1$ and $K=3$ disagree ---
the decision depends on the flexibility of the classifier.)

\medskip
\textbf{(c) Effect of $K$.} Small $K$ (e.g.\ $K=1$): very flexible, jagged decision boundary
$\Rightarrow$ \emph{low bias, high variance}. Large $K$: averaging over many neighbours gives a
smooth, almost linear boundary $\Rightarrow$ \emph{high bias, low variance}. The training error
rate of 1-NN is $0$ (ignoring ties): each training point is its own nearest neighbour --- which
says nothing about its test error.

\medskip
\textbf{(d) Bayes classifier.} It assigns each $x$ to the class $j$ with the largest \emph{true}
conditional probability $\Pr(Y=j\mid X=x)$; it minimises the expected test error rate, and its
error (the \emph{Bayes error rate}) is the lower bound for all classifiers. It is unattainable
in practice because the true conditional distribution of $Y$ given $X$ is unknown. KNN mimics
it by \emph{estimating} these conditional probabilities locally from the $K$ nearest training
observations and then classifying to the largest estimated probability.
\end{solutionbox}
\fi

% ============================================================
\problem{4: Simple linear regression by hand}{24}

A caf\'e chain investigates the effect of local radio advertising on sales.
For $n=5$ weeks it records the number of radio spots per week $x$ and the weekly sales $y$
(in 1{,}000 EUR):

\begin{center}
\begin{tabular}{lccccc}
\toprule
Week $i$ & 1 & 2 & 3 & 4 & 5\\
\midrule
Spots $x_i$ & 1 & 2 & 3 & 4 & 5\\
Sales $y_i$ & 6 & 7 & 11 & 13 & 18\\
\bottomrule
\end{tabular}
\end{center}

You may use without proof: $\sum x_i = 15$, $\sum y_i = 55$, $\sum x_i^2 = 55$,
$\sum x_i y_i = 195$. Consider the model $Y=\beta_0+\beta_1 X+\varepsilon$.

\begin{enumerate}[label=(\alph*),leftmargin=7mm,itemsep=2.5mm]
  \item \pp{6} Compute the least-squares estimates $\hat{\beta}_1$ and $\hat{\beta}_0$, and
    write down the fitted regression line. Interpret both estimates in the context of the
    problem (with units).
  \item \pp{4} Compute the fitted values $\hat{y}_i$ and the residuals $e_i$ for all five
    weeks, and verify that $\sum_{i=1}^{5} e_i = 0$.
  \item \pp{3} Compute the residual sum of squares (RSS) and the residual standard error
    (RSE), and interpret the RSE.
  \item \pp{6} Compute $\operatorname{SE}(\hat{\beta}_1)=\operatorname{RSE}/\sqrt{\sum_i (x_i-\bar{x})^2}$
    and test $H_0\!:\beta_1=0$ against $H_1\!:\beta_1\neq 0$ at level $\alpha=0.05$ using the
    $t$-statistic. Also give the 95\,\% confidence interval for $\beta_1$.
    \emph{Use $t_{0.975,\,3}=3.182$ (97.5\,\% quantile of the $t$-distribution with 3 degrees
    of freedom).} State your conclusion in one sentence.
  \item \pp{3} Compute the total sum of squares (TSS) and the coefficient of determination
    $R^2$, and interpret $R^2$.
  \item \pp{2} Predict the weekly sales for a week with $x=6$ radio spots. Name one reason
    why this particular prediction should be treated with caution.
\end{enumerate}

\ifdefined\withsolutions
\begin{solutionbox}
\textbf{(a) Coefficient estimates.} $\bar{x}=15/5=3$, $\bar{y}=55/5=11$.
\[
S_{xy}=\sum x_iy_i-n\bar{x}\bar{y}=195-5\cdot 3\cdot 11=30,\qquad
S_{xx}=\sum x_i^2-n\bar{x}^2=55-5\cdot 9=10.
\]
\[
\hat{\beta}_1=\frac{S_{xy}}{S_{xx}}=\frac{30}{10}=3.000,\qquad
\hat{\beta}_0=\bar{y}-\hat{\beta}_1\bar{x}=11-3\cdot 3=2.000.
\]
Fitted line: $\hat{y}=2.000+3.000\,x$.
\emph{Interpretation:} each additional radio spot per week is associated with an expected
increase in weekly sales of $3.000$ thousand EUR ($=3{,}000$ EUR). The intercept $2.000$ is the
expected weekly sales ($2{,}000$ EUR) in a week without any spots --- a mild extrapolation,
since $x=0$ was not observed ($x$ ranges from 1 to 5).

\medskip
\textbf{(b) Fitted values and residuals.} $\hat{y}_i=2+3x_i$ and $e_i=y_i-\hat{y}_i$:
\begin{center}
\begin{tabular}{lccccc}
\toprule
$i$ & 1 & 2 & 3 & 4 & 5\\
\midrule
$\hat{y}_i$ & 5 & 8 & 11 & 14 & 17\\
$e_i$ & $1$ & $-1$ & $0$ & $-1$ & $1$\\
\bottomrule
\end{tabular}
\end{center}
$\sum_i e_i = 1-1+0-1+1 = 0$ \checkmark\ (this always holds for least squares with an intercept).

\medskip
\textbf{(c) RSS and RSE.}
$\operatorname{RSS}=\sum_i e_i^2 = 1+1+0+1+1 = 4.000$;
$\operatorname{RSE}=\sqrt{\operatorname{RSS}/(n-2)}=\sqrt{4/3}=1.155$ (thousand EUR).
\emph{Interpretation:} even if the true line were known, observed weekly sales would deviate
from it by about $1.155$ thousand EUR ($\approx 1{,}155$ EUR) on average --- an estimate of
$\sigma$, the standard deviation of $\varepsilon$.

\medskip
\textbf{(d) Standard error, $t$-test and confidence interval.}
\[
\operatorname{SE}(\hat{\beta}_1)=\frac{1.1547}{\sqrt{10}}=\frac{1.1547}{3.1623}=0.365,\qquad
t=\frac{\hat{\beta}_1-0}{\operatorname{SE}(\hat{\beta}_1)}=\frac{3.000}{0.3651}=8.216
\quad(\text{df}=n-2=3).
\]
Since $|t|=8.216 > 3.182=t_{0.975,3}$, reject $H_0$ at the 5\,\% level.
95\,\% confidence interval:
\[
3.000 \pm 3.182\cdot 0.3651 = 3.000 \pm 1.162 = [1.838,\; 4.162].
\]
Zero is not contained in the interval --- consistent with the test.
\emph{Conclusion:} there is a statistically significant (positive) linear association between
the number of radio spots and weekly sales.

\medskip
\textbf{(e) TSS and $R^2$.}
$\operatorname{TSS}=\sum_i (y_i-\bar{y})^2 = 25+16+0+4+49 = 94.000$;
\[
R^2=1-\frac{\operatorname{RSS}}{\operatorname{TSS}}=1-\frac{4}{94}=\frac{90}{94}=0.957.
\]
About $95.7\,\%$ of the variance of weekly sales is explained by the linear regression on the
number of radio spots. (In simple regression $R^2=r^2$, so $r=\sqrt{0.957}\approx 0.979$.)

\medskip
\textbf{(f) Prediction at $x=6$.}
$\hat{y}(6)=2+3\cdot 6=20.000$ thousand EUR ($=20{,}000$ EUR).
Caution: $x=6$ lies \emph{outside} the observed range $x\in[1,5]$ --- an extrapolation; the
linear relationship has not been observed there and may not hold. (Also, $n=5$ is very small.)
\end{solutionbox}
\fi

% ============================================================
\problem{5: Multiple linear regression --- interpreting Python output}{24}

A consultancy analyses the salaries of $n=60$ employees of a software company.
The variables are \texttt{salary} (annual salary in 1{,}000 EUR), \texttt{experience}
(professional experience in years) and the dummy variable \texttt{female}
($=1$ for female employees, $=0$ for male employees). The model
\[
\texttt{salary}=\beta_0+\beta_1\,\texttt{experience}+\beta_2\,\texttt{female}
+\beta_3\,(\texttt{experience}\times\texttt{female})+\varepsilon
\]
was fitted by least squares in Python (\texttt{statsmodels}), with the following output:

\begin{center}
\begin{tabular}{lrrrr}
\toprule
 & coef & std err & $t$ & $P>|t|$\\
\midrule
Intercept                    & 46.2000 & 1.100 & 42.000 & $<$0.0001\\
experience                   &  2.3500 & 0.250 &  9.400 & $<$0.0001\\
female                       & $-3.1000$ & 1.550 & $-2.000$ & 0.0504\\
experience:female            & $-0.8500$ & 0.350 & $-2.429$ & 0.0184\\
\bottomrule
\end{tabular}

\medskip
$n=60$; \quad $R^2 = 0.712$; \quad adjusted $R^2 = 0.697$;\\[0.5mm]
$F$-statistic $= 46.15$ on 3 and 56 degrees of freedom, \ Prob($F$-statistic) $< 0.0001$.
\end{center}

\begin{enumerate}[label=(\alph*),leftmargin=7mm,itemsep=2.5mm]
  \item \pp{6} Interpret all four estimated coefficients precisely in the context of the
    problem (with units), paying particular attention to the dummy variable and the
    interaction term.
  \item \pp{4} Which coefficients are significantly different from zero at level
    $\alpha=0.05$? Verify for the interaction term that the reported $t$-value is consistent
    with the reported coefficient and standard error. Should \texttt{female} be removed from
    the model because of its $p$-value? Discuss briefly.
  \item \pp{4} Write down the fitted regression equation separately for men and for women.
  \item \pp{4} Predict the annual salary of (i) a man with 10 years of experience and (ii) a
    woman with 10 years of experience. What is the estimated salary gap between women and men
    at 10 years of experience, and how can it be written in terms of $\hat{\beta}_2$ and
    $\hat{\beta}_3$?
  \item \pp{3} Interpret $R^2$. State the null hypothesis of the reported $F$-test and the
    conclusion that follows from the output.
  \item \pp{3} Model diagnostics:
    \begin{enumerate}[label=(\roman*),nosep,topsep=1mm]
      \item A plot of the residuals against the fitted values shows a clear U-shape
        (positive residuals for small and large fitted values, negative in between).
        What does this indicate, and name one possible remedy.
      \item In an extended model, the predictor \texttt{age} has a variance inflation factor
        of $\operatorname{VIF}=12$. What does this mean, and what is the practical consequence
        for the regression results?
    \end{enumerate}
\end{enumerate}

\ifdefined\withsolutions
\begin{solutionbox}
\textbf{(a) Interpretation of the coefficients} (salary in 1{,}000 EUR per year):
\begin{itemize}[nosep,leftmargin=5mm]
  \item \emph{Intercept} $46.20$: expected annual salary of a \emph{man}
    (\texttt{female}=0) with 0 years of experience: $46{,}200$ EUR.
  \item \emph{experience} $2.35$: for \emph{men}, each additional year of experience increases
    the expected salary by $2{,}350$ EUR. (Because of the interaction, this slope refers to the
    baseline group \texttt{female}=0 only.)
  \item \emph{female} $-3.10$: at 0 years of experience, women are expected to earn
    $3{,}100$ EUR \emph{less} than men --- the difference between the two intercepts.
  \item \emph{experience:female} $-0.85$: the experience slope for women is $0.85$ \emph{lower}
    than for men: women gain $2.35-0.85=1.50$ thousand EUR per additional year. Equivalently,
    the expected salary gap between men and women \emph{widens} by $850$ EUR with every year
    of experience.
\end{itemize}

\medskip
\textbf{(b) Significance at $\alpha=0.05$.}
Intercept and \texttt{experience}: $p<0.0001$ $\Rightarrow$ highly significant.
\texttt{experience:female}: $p=0.0184<0.05$ $\Rightarrow$ significant.
Consistency check: $t=\hat{\beta}_3/\operatorname{SE}(\hat{\beta}_3)=-0.85/0.35=-2.429$
\checkmark\ (matches the output; $p$-value from $t_{56}$).
\texttt{female}: $p=0.0504>0.05$ $\Rightarrow$ \emph{not} significant at exactly the 5\,\%
level, though only marginally. It should nevertheless \emph{not} be removed: by the
\emph{hierarchy principle}, if an interaction ($\texttt{experience}\times\texttt{female}$) is
kept in the model, the corresponding main effects should be kept as well --- otherwise the
interaction changes its meaning. Moreover, $p=0.0504$ is borderline evidence, not evidence of
no effect.

\medskip
\textbf{(c) Fitted equations per group.}
\begin{align*}
\text{Men } (\texttt{female}=0):\quad & \widehat{\texttt{salary}} = 46.20 + 2.35\,\texttt{experience}\\
\text{Women } (\texttt{female}=1):\quad & \widehat{\texttt{salary}} = (46.20-3.10) + (2.35-0.85)\,\texttt{experience}\\
 & \phantom{\widehat{\texttt{salary}}} = 43.10 + 1.50\,\texttt{experience}
\end{align*}
Two different intercepts \emph{and} two different slopes (that is exactly what the interaction allows).

\medskip
\textbf{(d) Predictions at 10 years of experience.}
\begin{align*}
\text{Man: } & 46.20 + 2.35\cdot 10 = 69.70 \;\Rightarrow\; 69{,}700 \text{ EUR}\\
\text{Woman: } & 43.10 + 1.50\cdot 10 = 58.10 \;\Rightarrow\; 58{,}100 \text{ EUR}
\end{align*}
Estimated gap: $58.10-69.70=-11.60$, i.e.\ women earn $11{,}600$ EUR less. In terms of the
coefficients: $\hat{\beta}_2+10\,\hat{\beta}_3=-3.10-8.50=-11.60$ thousand EUR.

\medskip
\textbf{(e) $R^2$ and $F$-test.}
$R^2=0.712$: the three model terms explain $71.2\,\%$ of the variance of salaries in the sample.
$F$-test: $H_0\!: \beta_1=\beta_2=\beta_3=0$ (none of the predictors is related to salary)
against the alternative that at least one $\beta_j\neq 0$. With $F=46.15$ (3 and 56 df) and
$p<0.0001$, $H_0$ is clearly rejected: the model is useful overall.

\medskip
\textbf{(f) Diagnostics.}
(i) A U-shaped residual-vs-fitted plot indicates \emph{non-linearity}: the linear model misses
systematic curvature in the relationship. Remedy: include a non-linear term, e.g.\ add
$\texttt{experience}^2$ (polynomial regression), or transform variables
(e.g.\ $\log(\texttt{salary})$).
(ii) $\operatorname{VIF}_j = 1/(1-R^2_{X_j\mid X_{-j}})$; $\operatorname{VIF}=12$ means that
regressing \texttt{age} on the other predictors gives $R^2\approx 0.917$: \texttt{age} is
almost a linear combination of the other predictors (strong \emph{multicollinearity}; common
rules of thumb flag VIF $>5$ or $>10$). Consequence: $\operatorname{SE}(\hat{\beta}_{\texttt{age}})$
is inflated by a factor $\sqrt{12}\approx 3.5$ relative to the uncorrelated case --- the
coefficient is estimated very imprecisely, its sign may flip between samples, and its $t$-test
loses power, even though the model's overall fit and predictions may be unaffected.
\end{solutionbox}
\fi

\vspace{6mm}
\begin{center}
\rule{0.35\textwidth}{0.4pt}\\[1mm]
\textit{End of the examination --- good luck!}
\end{center}

\end{document}
